\chapter{Practical Component Construction}
\label{ch:practical-components}
\chapterquote{A useful widget is not a clever subclass. It is a small contract whose state, commands, focus, and errors can be reviewed.}{The component construction rule}

\begin{chaptermap}
Core question & How should business applications build compound Swing components that look integrated, remain testable, and do not hide state in listeners?\\
Ingredients & Components, Actions, input maps, popup menus, text models, table headers, temporal values, MigLayout, UI defaults, accessibility, validation, lifecycle scopes, and screenshot tests.\\
Plain Swing baseline & Ordinary Swing is enough when the component names its model, owns its commands, respects the Look-and-Feel, and exposes reviewable state.\\
Corusco connection & Corusco contributes typed keys, command metadata, field models, table descriptors, problem targets, resources, help topics, and lifecycle scopes. The widget still remains ordinary Swing code.\\
Companion stance & This chapter uses focused sketches rather than a new runnable companion suite. Each sketch is meant to show the construction contract, not to be copied as a finished library.\\
\end{chaptermap}

\section{Component Construction Rules}

\begin{roadmapbox}
This opening section establishes the review contract used by every later
component. It starts from the value and command model, then works outward to
focus, keyboard access, accessibility, Look-and-Feel integration, lifecycle
cleanup, and testing. The goal is to make each later widget read as a small
system with visible state instead of as a clever subclass.
\end{roadmapbox}

Practical Swing applications eventually need controls that the standard widget set does
not provide directly. A toolbar needs a split button whose main press runs the common
command and whose arrow opens less common choices. A form needs a text field with a
browse button, a clear button, or a help button attached to the same visual unit. A
measurement field needs the typed number and unit to feel like one value, not two
unrelated widgets. A large table needs headers that sort, explain, filter, and expose
column visibility. A scheduling screen needs a calendar, a date-time chooser, or a span
selector whose model knows the difference between draft selection and committed value.

These controls are not exotic. They are ordinary business application furniture. They
become difficult only when they are treated as decorative subclasses instead of small
contracts. A component contract answers a few questions before it answers how pixels are
painted. What value does the user edit? What command does the primary gesture invoke?
What secondary gestures exist? What state disables the command? What object owns
validation problems? How does keyboard focus enter and leave the control? What should
assistive technology call it? What listener or popup must be removed when the screen
closes? If these answers are clear, the implementation can be plain.

\termdef{Compound component}{in this chapter} means a visible control built
from several Swing objects but intended to behave as one application control. It may be
a \api{JPanel} containing a text field and buttons. It may be a
\api{JButton} paired with a \api{JPopupMenu}. It may be a custom
\api{JComponent} that paints a calendar grid. It may be a header renderer and
mouse adapter installed on a
\api{JTableHeader}. The compound
nature is an implementation detail. The application should see a coherent value, command
surface, and lifecycle boundary.

\termdef{Widget model}{in this chapter} is the small state object that lets a
compound component avoid inventing state from current pixels. A split button may observe
a primary \api{Action} and a secondary action list. A unit field may expose raw text,
parsed amount, selected unit, and problem state. A calendar may expose visible month,
focused day, selected draft day, and committed value. A column header may observe table
sort state, descriptor metadata, and column visibility. The model may be an explicit
class, a record plus observable values, a Corusco field model, or a presenter-owned
object. The important point is that durable state has a home outside event handlers.

\begin{principlebox}{Model Before Decoration}
Start every compound control by naming the value, commands, selection, and problem
state. Painting, borders, and popup placement come after the state contract is visible.
\end{principlebox}

The temptation is to begin with the visual trick. A developer subclasses
\api{JTextField} to draw a button at the right edge. Another developer adds a
mouse listener that opens a popup when the click lands near the arrow. A third developer
stores the chosen unit in a client property because it is convenient. The prototype
looks right. Then keyboard operation is added, validation wants to target the unit,
localization makes the unit label wider, dark mode changes border colors, a screenshot
test runs with larger fonts, and the component no longer has a clear owner for its
facts. The fix is not to avoid custom controls. The fix is to keep the visual trick
subordinate to the contract.

Most business compound controls should be assembled from existing components before they
are hand-painted. A
\api{JPanel} with a text field and a
narrow button can respect focus traversal, accessibility, input methods, text editing,
high-DPI scaling, and UI defaults with much less code than a custom text editor. A
header popup can be a normal
\api{JPopupMenu}. A split button can
begin as a button and a separate arrow button inside a no-gap panel. Custom painting is
appropriate when the control is visually one surface, when there is no useful stock
component, or when performance and geometry require it. Even then, the implementation
should borrow standard input, focus, and action patterns.

\begin{tabularx}{0.96\linewidth}{@{}p{0.22\linewidth}X@{}}
\toprule
Question & Practical rule\\
\midrule
What is edited? & Name a value model before constructing child components.\\
What is invoked? & Use \api{Action} or command metadata for user operations, not anonymous listeners as the only owner.\\
How does focus work? & Decide whether the control is one focus stop or several coordinated focus stops.\\
How is it described? & Set accessible name and description from resources or descriptor metadata.\\
How is it styled? & Read UI defaults and component metrics before hard-coding colors, borders, or sizes.\\
How is it closed? & Return or register cleanup handles for listeners, popups, timers, and model subscriptions.\\
\bottomrule
\end{tabularx}

Focus policy is the first practical fork. Some compound controls should be one focus
stop. A split button is usually one command surface: Tab lands on it, Space invokes the
primary action, Down or Alt+Down opens the menu, and Escape closes the menu. A unit
field may be one logical field but two focusable regions if users must arrow into the
unit selector, or it may keep text focus and open the unit popup with a shortcut. A
date-time span selector is usually several focus stops because start date, start time,
end date, end time, and presets are distinct edit targets. The decision should be
deliberate because keyboard and accessibility behavior follow from it.

Input maps belong near the component contract. Mouse-only compound controls are not
finished controls. If a user can click an arrow, there should be a keyboard path to the
same popup. If a user can click a clear buddy button, there should usually be Escape or
a command binding when the field is focused. If a table header has a context menu,
Shift+F10 or the platform popup key should reach it. Swing already has
\api{InputMap} and \api{ActionMap};
use them instead of spreading key handling through key listeners.

\begin{lstlisting}[caption={A compound control installs named actions, not key listeners.}]
final class FieldCommands {
    final Action clear = new AbstractAction("clear") {
        @Override
        public void actionPerformed(ActionEvent event) {
            model.clearDraft();
        }
    };

    void install(JTextField field) {
        field.getInputMap().put(KeyStroke.getKeyStroke("ESCAPE"), "clear");
        field.getActionMap().put("clear", clear);
    }
}
\end{lstlisting}

The example is intentionally small. The important part is not Escape. It is the fact
that clearing the field is a named operation on the model. A clear buddy button, a
keyboard shortcut, and perhaps a popup command can all share the same action. The model
can decide whether clearing is allowed. Tests can invoke the action without
manufacturing a mouse event. A screen can close the binding if the action observes
external state.

Look-and-Feel respect is another place where prototypes drift. A component that looks
acceptable under one FlatLaf theme may look broken under a compact density, larger
default font, high-contrast theme, or platform Look-and-Feel. Prefer existing borders,
colors, insets, and component orientation. When a compound component wants to look like
one field, share the field's border and background rather than drawing a new unrelated
rectangle. When a popup button is attached to a field, size the button from font metrics
and UI defaults, not from a fixed pixel width chosen on one machine.

\begin{detailbox}{Client Properties Are Not a Design}
Client properties are useful for small hints to UI delegates and local installers. They
are a poor place for durable application state such as the selected unit, the current
filter, or a validation target. If another class must understand the value, give the
value a real model.
\end{detailbox}

Accessibility should be planned when the control is designed, not added by guessing
after screenshots pass. A buddy button whose accessible name is "button" is unfinished.
A status badge that communicates only by color is unfinished. A calendar day cell that
is painted but not keyboard reachable is unfinished. The exact depth of accessibility
support depends on the control, but every business widget should expose a name, a useful
description where needed, keyboard operation, focus indication, and non-color ways to
read validation or state.

Testing follows the same contract. Model tests prove parsing, formatting, selection, and
validation independent of Swing. Swing tests prove that actions are installed, focus
traversal is reachable, popups open from keyboard and mouse paths, disabled state is
reflected, and cleanup removes listeners. Screenshot tests prove density and
Look-and-Feel integration. Do not make the screenshot test carry the entire design. A
screenshot can show that a unit field looks integrated; it cannot prove that the unit is
part of the committed value.

\begin{principlebox}{Lessons Learned From Component Construction Rules}
The stable part of a compound component is not the child-component tree. It is
the contract that says what value is edited, what commands can run, what problem
state exists, how focus moves, and who closes resources. Once those facts are
explicit, Swing's ordinary mechanisms become enough for many controls that look
custom on the screen.
\end{principlebox}

\section{Corusco as the Construction Surface}

\begin{roadmapbox}
This section explains where Corusco belongs in practical component construction.
It does not replace Swing widgets. It supplies typed identity, generated
resources, command descriptors, field models, problem targets, table
descriptors, and lifecycle scopes. The section shows how those facts reduce the
amount of state that a custom Swing component must invent.
\end{roadmapbox}

Corusco is useful in this chapter precisely because these controls sit between low-level
Swing and application meaning. A split button is a Swing surface, but its commands are
application commands. A unit field is a Swing surface, but its value is a form field
with parsing and validation. A table header is a Swing surface, but its identity comes
from table descriptors. A temporal selector is a Swing surface, but its commit policy
belongs to a field model and workflow. Corusco gives these middle layers names so the
component does not need to smuggle business facts through labels, client properties, or
anonymous listeners.

The right use of Corusco is not to generate every custom widget. The right use is to let
generated and typed facts feed ordinary widgets. Generated field keys can name the unit
field. Generated command descriptors can create the split button's primary action and
popup actions. Generated table descriptors can name the header column under the pointer.
Generated problem codes can target the amount portion of a measurement. A lifecycle
scope can own the listener that refreshes a popup menu. The custom Swing code becomes
smaller because it no longer invents application identity.

\begin{principlebox}{Corusco Keeps Widgets Honest}
A practical component should consume Corusco keys, commands, resources, problems,
descriptors, and scopes as inputs. It should not recreate those facts from component
text or layout position.
\end{principlebox}

The useful boundary is an adapter. A reusable component may expose a small plain Swing
interface such as
\api{MeasurementFieldView}. A Corusco
binding then adapts a generated field model to that interface. The component does not
need to know about annotation processing. The generated form model does not need to know
about every custom component class. The adapter is the place where the two contracts
meet. This is the same design that makes renderers, editors, and table descriptors
manageable: each side has a stable interface.

\begin{lstlisting}[caption={A custom field adapter connects Corusco state to Swing.}]
final class MeasurementFieldBinding implements Binding {
    private final MeasurementFieldView view;
    private final MeasurementFieldModel field;
    private final Subscription subscription;

    MeasurementFieldBinding(
            MeasurementFieldView view,
            MeasurementFieldModel field) {
        this.view = view;
        this.field = field;
        this.subscription = field.changes().subscribe(this::refresh);

        view.addChangeListener(event -> {
            field.setRawAmount(view.rawAmount());
            field.setUnit(view.unit());
        });
        refresh(field.snapshot());
    }

    private void refresh(MeasurementSnapshot snapshot) {
        view.setRawAmount(snapshot.rawAmount());
        view.setUnit(snapshot.unit());
        view.setProblems(snapshot.problems());
    }

    @Override
    public void close() {
        subscription.close();
    }
}
\end{lstlisting}

The class names are illustrative, but the pattern should be familiar from the binding
chapters. The view exposes widget facts. The field model owns editing facts. The binding
transfers changes and closes subscriptions. That closure is not ceremonial. Without it,
an old field can keep receiving updates after a dialog closes or after a card is
replaced. Corusco scopes make that cleanup part of the construction vocabulary rather
than a memory of which listener was installed where.

Command metadata is equally valuable. Split buttons and command bars often grow
inconsistent because one developer sets button text, another sets menu text, another
sets a tooltip, and a fourth wires a keyboard shortcut. A Corusco command descriptor
lets the same command identity feed all surfaces. The custom split button sees
\api{Action} objects, but those
actions can be created from typed command keys and resource-backed metadata.

\begin{lstlisting}[caption={Generated command metadata can feed a split button.}]
SplitButtonModel exportCommands =
        new SplitButtonModel(
                actions.action(CustomerActions.EXPORT_PDF),
                List.of(
                        actions.action(CustomerActions.EXPORT_CSV),
                        actions.action(CustomerActions.EXPORT_XLSX),
                        actions.action(CustomerActions.PRINT_PREVIEW)),
                resources.text(CustomerResources.EXPORT_MORE));

SplitButton export = new SplitButton(exportCommands);
scope.add(exportCommandsBinding(export, presenter.exportState()));
\end{lstlisting}

The sketch assumes generated keys such as \api{CustomerActions.EXPORT\_PDF} and resource
keys such as
\api{CustomerResources.EXPORT\_MORE}.
The exact names in a codebase can differ. The point is that text, help, enablement, and
identity are not copied by hand into the split button. The button can remain a reusable
Swing class. The screen decides which generated commands belong to that instance.

Field keys give custom controls the same recovery path as stock controls. A problem
summary can route to a generated field key whether the field is a plain
\api{JTextField}, a buddy field, a measurement field, or a date-time span selector. The
binding knows how to focus the right component. The problem model knows which field is
invalid. The summary does not need to know the child component tree. This is especially
important for compound controls whose visible error state may cover several children.

\begin{lstlisting}[caption={A problem target should name the field contract.}]
Problem amountProblem =
        Problem.parse(
                OrderFields.SHIPPING_WEIGHT,
                OrderProblems.SHIPPING_WEIGHT_AMOUNT_PARSE);

Problem unitProblem =
        Problem.invalidChoice(
                OrderFields.SHIPPING_WEIGHT,
                OrderProblems.SHIPPING_WEIGHT_UNIT_UNAVAILABLE);
\end{lstlisting}

The two problems target the same logical field and different problem codes. A UI policy
may show one border around the whole measurement field and a tooltip that distinguishes
amount and unit. A form summary may navigate to the measurement field. A test may assert
the generated problem code. None of those behaviors require scraping the text "kg" from
the unit button.

Resources matter more for custom widgets because custom widgets are where hard-coded
strings often sneak in. Button arrows, popup titles, disabled reasons, day names, unit
abbreviations, quick range names, accessible descriptions, and validation summaries all
need product language. Corusco resource keys give the component installer a typed source
for those strings. The reusable component should accept strings, actions, or metadata;
the screen should resolve them through resources.

\begin{tabularx}{0.96\linewidth}{@{}p{0.28\linewidth}X@{}}
\toprule
Widget fact & Corusco source\\
\midrule
Split button command id & Generated action key or command descriptor.\\
Buddy field target & Generated field key and component key.\\
Measurement unit problem & Generated field key plus generated or handwritten problem code.\\
Header menu column & Generated or handwritten table descriptor column key.\\
Calendar value & Field model with typed date or date-time value.\\
Popup listener lifetime & View scope or detachable scope.\\
Accessible text & Resource key, help topic, or descriptor metadata.\\
\bottomrule
\end{tabularx}

This table is the mental model for the rest of the chapter. The widget is not asked to
become Corusco-specific. It is asked to have enough named seams for Corusco-generated
facts to enter. A split button constructor that accepts only literal strings is hard to
use well. A constructor that accepts actions and a popup name is easy to feed from
command metadata. A unit field that exposes only child components is hard to bind. A
unit field that exposes raw amount, unit, and problems is easy to adapt.

There is also a migration benefit. A legacy screen may have custom controls already.
Corusco can be introduced by wrapping those controls with adapters instead of replacing
them. First, define field keys and command keys. Then write bindings that connect
existing widgets to generated form models and command descriptors. Then move validation
into problem models. Finally, improve the visual component if needed. The screen becomes
more maintainable without a disruptive rewrite of every widget at once.

\begin{figure}[htbp]
\centering
\includegraphics[width=0.96\linewidth]{\practicalSplitButtonsFigure}
\caption{Split buttons are command surfaces. The primary button, popup menu, resources, disabled state, and cleanup scope can all be fed from Corusco command metadata while the widget remains ordinary Swing.}
\label{fig:practical-split-buttons}
\end{figure}

\begin{principlebox}{Lessons Learned From the Corusco Surface}
Corusco is most helpful at the boundary between application meaning and Swing
mechanics. A custom component should consume generated facts through small
adapters, not know how the whole application is generated. That keeps the widget
ordinary, testable, and reusable while preserving typed field and command
identity.
\end{principlebox}

\section{Split Buttons and Popup Buttons}

\begin{roadmapbox}
This section builds the command model for a split button. It separates the
primary action from secondary menu actions, chooses mouse and keyboard
semantics, discusses disabled reasons and popup lifetime, and shows why a menu
should be rebuilt from command state rather than treated as static decoration.
\end{roadmapbox}

A split button has two related surfaces. The main surface invokes the default operation.
The secondary surface opens a menu of related operations. A popup button has one surface
whose purpose is to open a menu. Both controls appear small, but they expose several
contracts: primary action, secondary actions, keyboard access, popup placement, disabled
state, accessible names, and command identity. Most bugs appear when those contracts are
distributed across anonymous mouse listeners.

The first design decision is whether the primary operation is real. If there is a clear
default, use a split button: Save runs Save, and the arrow offers Save As, Save All, or
Save Copy. Export runs the last chosen export format, and the arrow offers PDF, CSV,
XLSX, and print preview. If there is no default, use a popup button: Add opens a menu of
Add Customer, Add Contact, Add Address, and Add Task. Do not create a split button whose
primary half does nothing or opens the same menu. Users learn quickly whether a command
surface has a main action.

\begin{tabularx}{0.96\linewidth}{@{}p{0.25\linewidth}X@{}}
\toprule
Shape & Use when\\
\midrule
Split button & One operation is the normal command and alternatives are closely related.\\
Popup button & The control chooses among operations and no operation deserves primary status.\\
Menu button in toolbar & Space is tight and the command group is secondary to the screen.\\
Plain button plus nearby menu & The menu is not part of the same semantic command.\\
\bottomrule
\end{tabularx}

The primary action should be an
\api{Action}. The menu items should
also be
\api{Action} instances or commands adapted to actions. That gives text,
enabled state, icons, accelerator descriptions, selected state, and disabled reasons a
common home. Corusco command metadata can feed those actions, but the Swing control
should not care whether the metadata was generated or handwritten. The button observes
actions. It does not invent command policy.

\begin{lstlisting}[caption={A split button model separates primary and menu actions.}]
record SplitButtonModel(
        Action primary,
        List<Action> secondary,
        String popupActionName) {

    boolean hasPopup() {
        return !secondary.isEmpty();
    }
}
\end{lstlisting}

This tiny record is enough to prevent a common failure. Without a model, code often
constructs the main button in one method and the popup menu in another, with no obvious
statement that the actions belong to one command group. With a model, the component can
update its enabled state when the primary action changes, rebuild or refresh menu items
from the secondary list, and expose an accessible description that names the menu. A
richer implementation may observe an observable list of commands, but the shape is the
same.

There are two implementation styles. The conservative style uses a
\api{JPanel} with two buttons: one normal button and one narrow arrow button.
The panel removes the gap between them and shares a height. This approach respects
button UIs, focus painting, and accessibility with little custom code. The tighter style
subclasses
\api{JButton} or creates a custom
\api{JComponent} that paints the main and arrow zones as one surface. That
approach can look better but must reimplement hit testing, focus, keyboard, pressed
state, rollover, disabled state, and accessible actions. Start with two buttons unless
the product genuinely needs a one-piece visual.

\begin{lstlisting}[caption={A conservative split button assembled from two buttons.}]
final class SplitButton extends JPanel {
    private final JButton main = new JButton();
    private final JButton arrow = new JButton("v");
    private final JPopupMenu menu = new JPopupMenu();

    SplitButton(SplitButtonModel model) {
        super(new MigLayout("insets 0, gap 0", "[]0[24!]", "[fill]"));
        main.setAction(model.primary());
        arrow.setFocusable(false);
        arrow.setToolTipText(model.popupActionName());

        for (Action action : model.secondary()) {
            menu.add(new JMenuItem(action));
        }

        arrow.addActionListener(event -> showMenu());
        add(main, "growy");
        add(arrow, "growy");
    }

    private void showMenu() {
        menu.show(this, 0, getHeight());
    }
}
\end{lstlisting}

This is not a finished library class. It omits resource lookup, component orientation,
listener cleanup, disabled reason composition, and keyboard installation. It shows the
baseline: two child buttons, one popup, actions as the command surface. The arrow button
is not focusable because focus remains on the main command group. In a screen where the
arrow must be independently reachable, make it focusable and describe that traversal in
the component contract.

Popup placement should be boring. Show the menu below the control for normal
left-to-right horizontal toolbars. Respect component orientation when aligning the
popup. If the button is near the bottom of the screen, \api{JPopupMenu} and the platform
may adjust placement. Do not manually clamp coordinates unless a real bug requires it.
The menu belongs to the component lifecycle: items may be rebuilt when commands change,
but the popup object itself should not hold stale references after the owning screen
closes.

Keyboard behavior deserves explicit names. Space and Enter invoke the primary action
when the split button is focused. Down, Alt+Down, or a named menu key opens the popup.
Escape closes the popup. Shift+F10 or the popup key may open the menu as a context
command. These bindings should live in the split button or installer, not in the
containing toolbar. The containing screen should not need to know which key opens the
arrow zone.

\begin{lstlisting}[caption={Keyboard bindings make the popup a command path.}]
void installPopupKeys(JComponent control, Action openPopup) {
    InputMap input = control.getInputMap(JComponent.WHEN_FOCUSED);
    ActionMap actions = control.getActionMap();

    input.put(KeyStroke.getKeyStroke("alt DOWN"), "openPopup");
    input.put(KeyStroke.getKeyStroke("F10"), "openPopup");
    actions.put("openPopup", openPopup);
}
\end{lstlisting}

The F10 example may be adjusted for platform conventions. The point is that the popup
path is a named action. It can check whether any menu items are enabled, update menu
state before showing, and log diagnostics if command metadata is missing. A key listener
would make those checks harder to share with the arrow button.

Disabled state is subtle. If the primary action is disabled but at least one secondary
action is enabled, should the split button be enabled? Many products choose to keep the
arrow available while disabling only the main surface. A two-button implementation
handles this naturally: the main button follows the primary action, and the arrow
follows whether the menu has enabled entries. A single-surface implementation must
represent a partially enabled state. The decision should be documented. Otherwise a user
may lose access to secondary commands because the default command is disabled.

\begin{pitfallbox}{One Enabled Flag}
Do not give a split button only one enabled flag unless the primary and secondary
surfaces truly share the same availability. A disabled default command does not always
mean the menu is useless.
\end{pitfallbox}

The popup menu should be refreshed at the moment it opens if command state can change.
Refreshing does not mean rebuilding every item from scratch on every mouse move. It
means updating enabled state, selected state, text, and perhaps visibility from the
command model. If commands are expensive to compute, the presenter should prepare them
before the menu opens. The component should not perform service calls or domain queries
during popup display.

Accessibility for split buttons should expose both operations. A simple two button
implementation gets two accessible children, but that may not be the best user
experience if the visual control is one unit. A wrapper can set an accessible name such
as "Export" and an accessible description such as "Press Space to export using the
default format. Press Alt+Down for export formats." The exact wording belongs in
resources. The important part is that the secondary action is discoverable without
seeing the arrow.

Review split buttons by walking one command. Find the command id. Find its resource
text. Find the primary action. Find the popup action. Find the key binding. Find
disabled state. Find the cleanup path. If any of these facts exist only in a mouse
listener, the control is probably still a prototype.

\begin{principlebox}{Lessons Learned From Split Buttons}
A split button is a command cluster, not a button with a decorative arrow. The
primary action, secondary actions, popup origin, default selection, disabled
state, and keyboard path all need one command model. Corusco command metadata
keeps labels, help, permissions, and disabled reasons consistent across the main
button, menu items, and tests.
\end{principlebox}

\section{Design walk: decisions behind a split button}

\begin{roadmapbox}
This design walk slows down the implementation decision process. It asks which
action is primary, how defaults are chosen, what state can disable individual
commands, how popup content is refreshed, and how the component can stay
reviewable when the command set changes at runtime.
\end{roadmapbox}

The first real decision is not how to draw the arrow. It is what the primary operation
means. A split button should have a stable default action that a user can invoke without
reading the menu. If the application cannot name that default confidently, the component
should be a popup button instead. This decision affects tests, resources, help,
telemetry, and keyboard behavior. "Export" may be a popup-only command if the user must
always choose a format. "Export PDF" may be the primary action if PDF is the product
default and the menu contains less frequent alternatives. The code should express that
decision by giving the primary action its own key, not by making the first menu item
special.

The second decision is whether secondary commands are static or contextual. Static
commands are known at construction time: Export CSV, Export XLSX, Print Preview.
Contextual commands depend on selection, permissions, current row type, license, or
backend capability. A reusable split button should not compute those facts. It should
accept a model or command list whose state can be refreshed before the menu opens. The
presenter decides which commands exist. The component decides how to present them.

\begin{lstlisting}[caption={A contextual split-button model exposes command availability.}]
record CommandChoice(
        Action action,
        boolean visible,
        Optional<String> disabledReason) {
}

interface SplitCommandSource {
    Action primaryAction();
    List<CommandChoice> secondaryChoices();
    void refreshForCurrentContext();
}
\end{lstlisting}

This shape prevents the menu from becoming a mini presenter. Before showing the popup,
the split button calls
\api{refreshForCurrentContext}. It
then adds visible choices and applies enabled state from their actions. Disabled reasons
can feed tooltips or accessible descriptions. If the current row does not allow CSV
export, the presenter says so. The component does not inspect the table selection or
domain object. That separation is what lets the same component appear in a toolbar,
command bar, or dialog footer.

Corusco command metadata makes the source smaller. A generated action key can resolve
text, icon, tooltip, help topic, accelerator description, and command id. The presenter
contributes only dynamic state: enabled, selected, visible, or disabled reason. The
Swing action adapter joins those facts. The split button receives ordinary \api{Action}
objects, but the action objects are not anonymous blobs. They carry generated identity.

\begin{lstlisting}[caption={A Corusco-backed command source keeps dynamic state small.}]
final class ExportCommandSource implements SplitCommandSource {
    private final CommandSet commands;

    ExportCommandSource(CommandSet commands) {
        this.commands = commands;
    }

    @Override
    public Action primaryAction() {
        return commands.swingAction(CustomerActions.EXPORT_PDF);
    }

    @Override
    public List<CommandChoice> secondaryChoices() {
        return List.of(
                choice(CustomerActions.EXPORT_CSV),
                choice(CustomerActions.EXPORT_XLSX),
                choice(CustomerActions.PRINT_PREVIEW));
    }
}
\end{lstlisting}

The component wiring is then a small state machine. In the idle state, the main surface
mirrors the primary action and the arrow surface mirrors whether there are visible
secondary choices. Pressing the main surface invokes the primary action. Pressing the
arrow surface refreshes choices, rebuilds or updates menu items, shows the popup, and
records that the popup is open. Escape closes the popup. A command invocation from the
popup closes the popup and then invokes the command. A model update while the popup is
open should update menu state if it is cheap, or close the popup if the command set has
changed structurally. The policy should be consistent.

The event path should be written down because it crosses Swing boundaries: mouse press
changes button model state, action invocation reaches command adapter, command adapter
calls presenter, presenter may start a task, task state updates command enablement,
binding refreshes action enabled state, and the split button repaints. The split button
should not reach into the task. It only observes action state. This is how Corusco keeps
a small visual control from becoming a workflow controller.

\begin{tabularx}{0.96\linewidth}{@{}p{0.27\linewidth}X@{}}
\toprule
Split-button decision & Implementation consequence\\
\midrule
Primary action exists & Space and Enter invoke it; tests assert the generated action key.\\
Secondary actions are contextual & Popup refreshes from a command source before showing.\\
Arrow remains available & Main and arrow enabled states are computed separately.\\
Menu items can be disabled & Disabled reason is displayed as tooltip or accessible description.\\
Commands are generated & Button, menu, shortcut text, help, and diagnostics share one id.\\
View can close & Popup listeners and command subscriptions are registered in a scope.\\
\bottomrule
\end{tabularx}

Failure modes follow from these decisions. If the primary command is disabled and the
arrow also disables, users may lose useful alternate commands. If the popup is rebuilt
from literal strings, resources and help drift. If menu items capture row indexes
instead of command context, sorting or selection changes can route an export to the
wrong row. If the popup listener is never removed, the menu may retain a presenter after
the screen closes. A split button is small enough that these bugs look disproportionate;
that is exactly why the contract should be explicit.

\begin{principlebox}{Lessons Learned From the Split-Button Walk}
The design pressure in a split button is the tension between visual compactness
and command clarity. The best implementation makes the compact surface honest:
the user can discover the default command, reach the secondary commands by
keyboard, read disabled reasons, and trust that the popup represents current
application state.
\end{principlebox}

\section{Text Fields With Buddy Buttons}

\begin{roadmapbox}
This section treats a buddy field as one form field with one or more attached
commands. It covers browse, search, clear, and help buttons; focus ownership;
validation decoration; sizing; keyboard shortcuts; and the adapter boundary
between a generated form field and a Swing panel that contains several child
components.
\end{roadmapbox}

A buddy button is a small command attached to a text field. Common examples are Browse,
Search, Clear, Reveal Password, Help, Pick From List, and Open Recent. The user
experiences one field with a nearby affordance. The code should treat the buddy as a
command related to the field model, not as an ornament glued to a text box.

\begin{figure}[htbp]
\centering
\includegraphics[width=0.96\linewidth]{\practicalBuddyUnitFigure}
\caption{Buddy fields and integral unit fields use ordinary Swing children, but the logical field is still named by Corusco field keys, problem codes, and binding scopes.}
\label{fig:practical-buddy-unit-fields}
\end{figure}

The implementation usually does not need a custom \api{JTextField}. A small
\api{JPanel} can contain a text field and one or more buttons, share a visual
border, and expose methods for reading the field and installing actions. MigLayout makes
this straightforward. A custom text field that draws icons inside itself should be
reserved for cases where the product requires the button to be inside the text rectangle
or where a Look-and-Feel already supports such decoration through client properties.

\begin{lstlisting}[caption={A buddy field wrapper keeps the text component visible.}]
final class BuddyTextField extends JPanel {
    private final JTextField text = new JTextField();
    private final JButton buddy = new JButton();

    BuddyTextField(Action buddyAction) {
        super(new MigLayout("insets 0, gap 0", "[grow,fill]0[28!]", "[fill]"));
        setBorder(UIManager.getBorder("TextField.border"));
        setBackground(text.getBackground());

        text.setBorder(BorderFactory.createEmptyBorder());
        buddy.setAction(buddyAction);
        buddy.setFocusable(false);

        add(text, "growx");
        add(buddy, "growy");
    }

    JTextField textComponent() {
        return text;
    }
}
\end{lstlisting}

The wrapper borrows the text field border and removes the child field border so the two
children look like one control. A production class would handle disabled state,
component orientation, high-DPI insets, accessible names, and Look-and-Feel updates. It
might also override \api{updateUI} to refresh the borrowed border and background when
the Look-and-Feel changes. The value still lives in the real \api{JTextField}, so input
methods, document filters, undo, selection, copy, paste, and accessibility remain
standard.

The buddy button's focusability is a design choice. A Clear button on a search field
often should not take focus; pressing it should clear the field and leave the insertion
caret ready for the next query. A Browse button that opens a file chooser may be a real
focus stop because keyboard users need to reach it separately. A Help button may be
reachable through F1 instead of Tab. Do not set \api{setFocusable(false)} mechanically.
Decide which commands are field shortcuts and which are independent controls.

\begin{tabularx}{0.96\linewidth}{@{}p{0.25\linewidth}X@{}}
\toprule
Buddy command & Focus and keyboard policy\\
\midrule
Clear & Usually not a focus stop; Escape or a named clear shortcut can invoke it while text retains focus.\\
Browse & Often a focus stop; Space opens the chooser, and the field remains editable directly.\\
Search & Often default command for the field; Enter invokes search while a button may also be visible.\\
Help & Usually F1 plus optional visible button; accessible description should name the target field.\\
Reveal password & May be press-and-hold or toggle; state must be visible and accessible.\\
\bottomrule
\end{tabularx}

Search fields show the difference between text state and command state. The text
document owns the current query draft. The presenter or field model owns whether the
query is valid, whether search is running, and whether clearing is available. The Search
action reads committed or current query according to policy. The Clear action clears the
draft and invalidates outstanding search results if the screen uses generation tokens.
The component should not start background work directly because that would mix visual
construction with workflow policy.

\begin{lstlisting}[caption={Search and clear actions share one field model.}]
final class SearchFieldModel {
    private String query = "";

    boolean canSearch() {
        return !query.isBlank();
    }

    void setQuery(String query) {
        this.query = query;
    }

    void clear() {
        setQuery("");
    }
}
\end{lstlisting}

In production this model may be observable and may store debounce state, normalization,
validation problems, and a generation counter. The sketch keeps only the value because
the ownership point matters more than the machinery. The component is free to show a
clear buddy only when \api{canSearch} is true, but the decision derives from model
state. It is not inferred from whether the button happens to be visible.

File and path fields add another concern: chooser results are committed through an
interaction that may be cancelled. The Browse action should read the current field value
as the initial directory or selected file, open the chooser, and write the result only
if the user approves. The raw text may remain invalid while the user types. The chooser
result may need conversion to a normalized path, display text, or domain object. Keep
these steps separate so validation can explain whether the problem is "path text cannot
be parsed", "file does not exist", "file exists but is not readable", or "selected item
is not allowed for this operation".

\begin{lstlisting}[caption={A browse action commits only an approved chooser result.}]
final class BrowsePathAction extends AbstractAction {
    private final JTextField field;
    private final JFileChooser chooser;

    BrowsePathAction(JTextField field, JFileChooser chooser) {
        super("Browse");
        this.field = field;
        this.chooser = chooser;
    }

    @Override
    public void actionPerformed(ActionEvent event) {
        chooser.setSelectedFile(new File(field.getText()));
        if (chooser.showOpenDialog(field) == JFileChooser.APPROVE_OPTION) {
            field.setText(chooser.getSelectedFile().getAbsolutePath());
        }
    }
}
\end{lstlisting}

The action is intentionally plain. A richer screen would route through a presenter,
resources, problem models, and perhaps recent locations. The important contract is that
cancellation leaves the field alone and approval writes through the same text path as
ordinary editing. The file chooser does not become a hidden second source of truth.

Layout details matter because buddy fields appear in dense forms. The wrapper should
expose a useful baseline so labels align with fields. If the wrapper is a \api{JPanel},
its baseline should normally delegate to the text field. If it cannot, MigLayout may
align the panel as a component but not the text inside it, producing small visual drift.
Baseline problems are especially visible when a form mixes plain fields, combo boxes,
buddy fields, and unit fields. Test with the real Look-and-Feel and at least one larger
font.

Validation decoration should attach to the field as a logical target, not to whichever
child happens to draw the border. If the wrapper owns the border, then error border or
problem marker installation should know that the wrapper is the visual field surface. If
the text child owns the border, decoration can remain on the child. Choose one. A common
failure is to decorate the inner text field while the wrapper draws the visible border,
causing errors to be invisible.

\begin{pitfallbox}{Two Borders, No Field}
If both the wrapper and the child text field draw normal borders, the control looks
heavy and validation decoration may land in the wrong place. Decide which component owns
the field surface.
\end{pitfallbox}

Buddy fields also need good accessible relationships. The visible label may label the
wrapper, the text field, or both. The buddy button should name its target: "Browse for
report folder", "Clear customer search", or "Show help for account number". A screen
reader should not encounter a row of unlabeled small buttons after a series of fields.
Resource-backed names solve this without hard-coding product text into the widget class.

Lifecycle is usually simple. Document listeners, action listeners, popup menus, chooser
listeners, delayed search timers, and validation subscriptions belong to the screen or
field binding that installed them. If the buddy field is reusable, expose an install
method that returns a cleanup handle or accept a scope into which listeners are
registered. A search field with a debounce timer is not just a component; it owns
scheduled work that must stop when the screen closes.

\begin{principlebox}{Lessons Learned From Buddy Fields}
Buddy buttons work when the field remains the value owner and the buttons remain
commands against that value. The visual border can make the control look like one
unit, but the model is what makes it behave like one unit. Clear, browse, search,
and help should all route through named actions that tests and Corusco bindings
can see.
\end{principlebox}

\section{Design walk: model and wiring for a buddy field}

\begin{roadmapbox}
This walk follows a buddy field from model to event wiring. It chooses the raw
text owner, decides how document changes become draft changes, attaches clear
and browse commands, places validation problems, and explains how the binding
avoids calling application services directly from the Swing view.
\end{roadmapbox}

The central decision for a buddy field is whether the buddy is an editing command or a
workflow command. Clear is usually an editing command: it changes the field draft.
Browse is partly editing and partly workflow: it opens a chooser and commits a chosen
path. Search is usually a workflow command: it submits the current query to a presenter.
Help is neither editing nor workflow; it routes to documentation for the field. Treating
all buddy buttons as "small buttons beside a field" hides these differences and leads to
inconsistent focus, validation, and testing.

The model should therefore name both the draft and the available field commands. A
clearable search field has raw query text, normalized query, problem state, dirty state,
pending search state, and commands. A path field has raw path text, parsed path, chooser
command, recent paths, and problems. A password field with reveal has raw secret text
and reveal state. The same wrapper component can present them, but the binding policy is
different.

\begin{lstlisting}[caption={A buddy field model names draft and field commands.}]
interface BuddyFieldModel {
    String rawText();
    void setRawText(String text);
    ProblemSet problems();
    CommandSet fieldCommands();
}
\end{lstlisting}

A Corusco binding can adapt generated field models to this interface. The generated
field key identifies the text field. A command key identifies Clear or Browse. A help
topic identifies F1 behavior. A problem set identifies validation. The Swing wrapper
only needs to expose child components and logical slots: text editor, primary buddy
button, optional secondary buddy button, visual field surface. The binding decides which
generated facts go into those slots.

Event wiring should be one-directional at each step. Document changes update the field
model. Field model changes refresh the text only when the change comes from outside the
document or when a reset is requested. The clear action calls the field model, not
\api{JTextField.setText} directly. The
browse action returns an approved value to the field model, not to an unrelated client
property. Validation observes the model and decorates the field surface. This avoids
feedback loops where document listeners respond to programmatic updates as if the user
typed them.

\begin{lstlisting}[caption={Buddy field wiring keeps commands behind the model.}]
Binding bindBuddyField(
        BuddyTextField view,
        TextFieldModel<String> field,
        Command clearCommand,
        BindingScope scope) {

    scope.add(SwingEditors.bindText(view.textComponent(), field));
    scope.add(CommandBehaviors.bindButton(view.buddyButton(), clearCommand));
    scope.add(problemDecoration(view.fieldSurface(), field.problems()));
    return scope;
}
\end{lstlisting}

The names are illustrative, but the ownership is precise. The editor binding owns
document synchronization. The command behavior owns button state and invocation. The
problem decoration owns visual feedback. The scope owns cleanup. Without these
boundaries, a developer reviewing the field must read through document listeners, action
listeners, and border changes to discover what the field actually does.

Focus policy should be written as a table before implementation. If Clear is not a focus
stop, then Tab moves from the field to the next form control, and Escape or a clear
shortcut invokes Clear. If Browse is a focus stop, Tab moves to Browse and Space opens
the chooser. If Help is available through F1, the visible help buddy may be
non-focusable or may be included only for discoverability. These choices are product
choices. Corusco component keys help tests assert the actual traversal: a test can find
the field and buddy by key rather than by text.

Chooser commands need a cancellation policy. Open chooser, user cancels, field remains
unchanged. Open chooser, user approves, value is normalized and written through the
model. Open chooser, chosen file fails validation, raw text may still be written but a
problem appears. Open chooser while a screen is closing, result should be ignored or
prevented. Those branches should not be hidden in the component. A presenter-owned
command can handle them and publish the resulting field state.

The failure modes are familiar. A clear button that calls
\api{setText("")}
directly may bypass dirty-state policy. A Browse button that writes a file object to a
client property may desynchronize display text and committed value. A Search button that
starts a repository call directly may outlive the screen. A non-focusable Browse button
may be unreachable from the keyboard. A validation border applied to the child field may
be invisible because the wrapper owns the visible border. The construction review should
explicitly ask for each of these paths.

\begin{principlebox}{Lessons Learned From the Buddy-Field Walk}
The important implementation choice is where side effects live. The field view
may open a chooser or request a search, but committed application changes should
flow through the presenter and model. That separation is what lets a generated
form, a validation layer, and an ordinary Swing control cooperate without hidden
listener state.
\end{principlebox}

\section{Text Fields With Integral Mini Combos}

\begin{roadmapbox}
This section designs an amount-plus-unit field that looks like one text field.
It chooses the model shape, explains why the unit is not a separate combo box,
handles parsing and formatting, positions the popup, preserves baseline and
border integration, and connects validation and form binding to both amount and
unit state.
\end{roadmapbox}

A measurement field such as \texttt{[ 120 kg v ]} looks like one text field. The number
and unit form one value. The unit selector is visually attached to the right side of the
field, but it is not a separate form column. Users read "120 kg" as one measurement.
Validation reads it as one field. Serialization reads it as one domain value. Treating
the unit as a separate combo box often creates alignment and ownership problems: the
label points at the number only, validation highlights one half, keyboard traversal
splits one value into two unrelated stops, and dirty state can change in either
component without a single field model noticing.

The value model should make the combined nature explicit. It usually contains raw text,
selected unit, parsed amount, and problem state. Raw text is needed because the user can
type an incomplete number. Selected unit is a constrained choice. Parsed amount may be
absent or invalid. The committed value should exist only when conversion and validation
succeed. A record such as
\api{MeasurementDraft} can express this separation.

\begin{lstlisting}[caption={A unit field model keeps raw text and selected unit together.}]
record MeasurementDraft(
        String rawAmount,
        MeasurementUnit unit,
        Optional<BigDecimal> parsedAmount,
        List<Problem> problems) {

    boolean canCommit() {
        return parsedAmount.isPresent() && problems.isEmpty();
    }
}
\end{lstlisting}

The committed domain value may be
\api{Measurement(BigDecimal amount,
MeasurementUnit unit)}. Do not force the draft to be the committed value. A user typing
a decimal may pass through states such as empty text, "-", "1.", or a locale-specific
decimal separator. The field should preserve those raw states while explaining problems.
The selected unit can still be valid while the amount is invalid, and changing the unit
may change validation because different units may have different ranges, precision, or
conversion policy.

The control can be assembled from a
\api{JTextField} and a small popup
button inside a shared border, similar to a buddy field. The difference is semantic: the
right-hand button is not an auxiliary command. It edits part of the same value. Its
visible text should be the current unit abbreviation, and the popup should present valid
units. If the unit text grows, the field layout must still hold. If the unit is disabled
or fixed by context, the selector may become a read-only suffix rather than a popup
button.

\begin{lstlisting}[caption={A unit field wrapper treats unit selection as field state.}]
final class MeasurementField extends JPanel {
    private final JTextField amount = new JTextField();
    private final JButton unitButton = new JButton();
    private final JPopupMenu units = new JPopupMenu();
    private MeasurementUnit unit;

    MeasurementField(List<MeasurementUnit> allowedUnits) {
        super(new MigLayout("insets 0, gap 0", "[grow,fill]0[]", "[fill]"));
        setBorder(UIManager.getBorder("TextField.border"));
        amount.setBorder(BorderFactory.createEmptyBorder());

        for (MeasurementUnit candidate : allowedUnits) {
            units.add(new JMenuItem(new SelectUnitAction(candidate)));
        }

        unitButton.addActionListener(event -> showUnits());
        add(amount, "growx");
        add(unitButton, "growy");
    }

    private void showUnits() {
        units.show(unitButton, 0, unitButton.getHeight());
    }
}
\end{lstlisting}

Again, the sketch is not a drop-in class. It omits observable model updates, resources,
component orientation, UI refresh, baseline delegation, and validation decoration. It
shows the shape: text and unit live in one wrapper. The unit button changes field state.
The popup is part of the field, not a screen-level menu.

The unit selector should have a clear keyboard path. Common choices are Alt+Down while
the amount field is focused, F4 by analogy with combo boxes, or Tab into a focusable
unit button. If the field is expected to be very fast for data entry, keeping text focus
and using Alt+Down is often better than creating a second Tab stop. If the unit changes
frequently and must be discoverable, a focusable unit segment may be better. The
decision should be consistent across the application. Do not make one measurement field
behave like a combo box and another behave like a button without a product reason.

Painting and borders are where unit fields usually become fragile. A separate
\api{JComboBox} beside a text field often has its own border, background, focus ring,
and arrow geometry. That can look like two fields. Using a button inside a shared
text-field border gives a more integrated appearance. The button should have reduced
insets and may use the field background. It still must show hover, pressed, disabled,
and focus state in a way that remains visible. Removing every visual state to make it
"clean" is not good design.

\begin{detailbox}{A Suffix Is Not Always a Button}
If the unit is fixed by context, render it as a non-focusable suffix label inside the
field wrapper. If the user can change it, make the selector keyboard reachable and
accessible. Do not hide an editable choice in a label.
\end{detailbox}

Baseline alignment should be tested with labels. The wrapper should align the amount
text baseline with neighboring text fields. The unit segment should not make the field
taller than a normal text field unless the Look-and-Feel does so consistently. If a unit
popup button uses a larger border or icon, forms become ragged. This is a good place to
rely on real font metrics and Look-and-Feel insets rather than fixed heights.

Validation should describe the whole measurement. A blank amount may be a required-field
problem. A nonnumeric amount is a parse problem. A value with too many decimals may be a
precision problem. A value outside allowed range may depend on the unit. A unit may be
disallowed for the current product or record. These problems should target the
measurement field, with optional subtargets for amount or unit when the UI can highlight
them separately. The presenter should not have to guess whether a problem belongs to the
text child or unit button by inspecting component coordinates.

\begin{lstlisting}[caption={Unit-aware validation returns field problems, not border choices.}]
List<Problem> validate(MeasurementDraft draft) {
    List<Problem> problems = new ArrayList<>();
    if (draft.rawAmount().isBlank()) {
        problems.add(Problem.required("weight"));
    }
    if (draft.parsedAmount().isEmpty()) {
        problems.add(Problem.parse("weight", "amount"));
    }
    if (!allowedUnits.contains(draft.unit())) {
        problems.add(Problem.invalidChoice("weight", "unit"));
    }
    return problems;
}
\end{lstlisting}

The strings in the sketch stand in for typed keys or generated problem codes. The
important point is that validation returns structured facts. The field binding decides
how to decorate: red border around the wrapper, marker near the amount, problem tooltip
on the unit segment, or a summary line below the form. Validation logic should not call
\api{setBorder}.

Formatting is a commit policy, not every-keystroke policy. Many applications format the
amount when the field loses focus or when the value is committed. Formatting on every
document change can fight the user, break undo, move the caret, and confuse input
methods. Unit changes may optionally reformat or convert the amount. The product must
decide whether changing from kg to lb converts the numeric value or merely changes the
unit attached to the same number. Both behaviors exist in real software. The field model
should make the policy explicit because silent conversion errors are expensive.

The field should not assume all units have short ASCII labels. Localized unit symbols
may be longer. Some products use business units such as "cases", "pallets", "hours",
"FTE", or "license seats". Some units require plural names in menus but abbreviations in
the field. The popup menu can display a long localized name while the field segment
displays a short label. Both should come from resources or unit metadata, not from enum
names.

Accessibility should describe the combined value. A screen reader should be able to
announce "Weight, edit text, 120 kilograms" and also expose how to change the unit. If
the unit segment is focusable, it needs an accessible name such as "Weight unit,
kilograms, popup button". If the unit is opened from the text field with Alt+Down, the
field's description should mention the shortcut. Validation summaries should identify
whether the amount or unit is wrong when that helps the user recover.

Form binding should treat the measurement as one field. A generated form model might not
know this exact custom component in v1, but the binding can adapt the component to the
same field concepts: raw input, parsed value, problems, touched state, dirty state, and
commit. The custom field should therefore expose methods or observable properties that a
binding can use. Do not require a binding to traverse child components and infer the
unit from button text.

\begin{lstlisting}[caption={A binding surface exposes measurement facts directly.}]
interface MeasurementFieldView {
    String rawAmount();
    void setRawAmount(String text);
    MeasurementUnit unit();
    void setUnit(MeasurementUnit unit);
    void setProblems(List<Problem> problems);
    void addChangeListener(ChangeListener listener);
}
\end{lstlisting}

This interface is deliberately minimal. A production API may use observable values and
typed keys. The key decision is that the custom component exposes domain-relevant facts
directly. It does not expose only \api{JTextField} and
\api{JButton} and hope every caller wires them consistently.

\begin{principlebox}{Lessons Learned From Integral Mini Combos}
The unit segment is part of the value, not an adjacent control that happens to
be nearby. Once the model contains raw amount, parsed amount, unit, format
policy, and problems, the view can be built from ordinary Swing pieces while
still feeling like one field. The hardest choices are conversion policy,
keyboard access, and problem targeting.
\end{principlebox}

\section{Design walk: amount, unit, conversion, and events}

\begin{roadmapbox}
This walk follows an amount field through the unpleasant details: invalid raw
text, unit changes after input, precision and rounding, canonical storage units,
popup selection, and commit events. It shows how the event path avoids silently
rewriting user text while still keeping typed domain values available.
\end{roadmapbox}

The mini-combo field has more decisions than its size suggests. The first decision is
whether changing the unit converts the amount or only changes the unit label. In a
weight field, changing 120 kg to lb may mean "convert 120 kg to 264.55 lb" or it may
mean "the user meant 120 lb". Both are legitimate in different products. The component
must not guess. The model needs a
\api{UnitChangePolicy}: preserve magnitude, convert physical quantity, clear
amount, or ask the presenter. The field's event path depends on that policy.

The second decision is what precision belongs to the amount. Kilograms may allow three
decimals, pieces may allow none, hours may allow quarter-hour increments, currency-like
units may have fixed scale. Precision is not only formatting; it affects parsing,
validation, caret behavior after format, and problem messages. The selected unit may
therefore contribute a converter and validator. A model that stores only
\api{BigDecimal} and enum unit is too
small for serious editing. It should also know raw text, parse state, unit metadata, and
formatting policy.

\begin{lstlisting}[caption={Unit metadata makes conversion and formatting explicit.}]
record UnitOption(
        MeasurementUnit unit,
        String abbreviationKey,
        String displayNameKey,
        int scale,
        UnitChangePolicy changePolicy) {
}
\end{lstlisting}

This metadata is a good Corusco resource boundary. The unit enum can remain stable.
Resource keys provide abbreviations and menu names. Validation rules use the unit
identity and scale. A unit popup displays localized names. The field suffix displays
abbreviations. Tests can use unit ids without depending on English labels. This keeps
the field maintainable when product language changes from "kg" to "kilograms" in a menu
but the field suffix remains short.

The event path has several branches. User types text: document binding writes raw text,
parser updates parse state, validators update problems, component repaints problem
state. User chooses unit: popup command changes selected unit, unit-change policy
decides whether amount text changes, parser and validators rerun, resources update
suffix text, component preserves or restores focus according to policy. Programmatic
reset: model replaces raw text and unit from committed value, document binding updates
field without marking user input as touched, problems clear or recompute. Each branch
should be testable without screenshots.

\begin{lstlisting}[caption={A unit change event is model work before Swing refresh.}]
void selectUnit(UnitOption option) {
    MeasurementDraft old = draft;
    MeasurementDraft next = unitChangePolicy.apply(old, option);
    draft = parseAndValidate(next);
    changes.emit(draft);
}
\end{lstlisting}

The component listens to \api{changes} and refreshes text, suffix, tooltip, and problem
marker. It does not implement conversion policy in the button listener. The button
listener says only "user selected this unit." That one sentence is the practical
difference between a maintainable widget and a widget that becomes impossible to reason
about after three product rules are added.

Focus and popup behavior are also policy. If the user opens the unit popup from the
keyboard while editing text, selecting a unit should usually return focus to the amount
field and preserve caret position where possible. If unit selection reformats the
amount, the caret may move to the end because the old position may no longer be
meaningful. If the unit segment is focusable, Tab order should make sense with
neighboring fields. If the field is read-only, the unit segment should not look
editable. These states need screenshot coverage because they are visual, but the rules
should be model tests first.

Validation can be layered. Parse problems belong to the amount text. Choice problems
belong to the unit. Cross-field problems belong to the measurement as a whole. Corusco
problem targets let the application keep one generated field key while using problem
codes or subtargets to explain recovery. The UI can draw one red outline around the
whole field and still show a message such as "kg is not allowed for this product" or
"amount has too many decimal places for pieces." A plain combo box beside a text field
often splits these problems into two unrelated form rows.

The deepest failure mode is silent unit drift. A form loads 120 kg. The user changes the
unit to lb. The UI shows 120 lb. The save path stores 120 kg because the binding never
observed the unit button. Or the reverse happens: the UI shows kg, but the model has lb.
This is why the binding surface should expose the combined measurement, not child
components. The Save command should read one field model, and that model should produce
one committed measurement or one structured problem set.

\begin{principlebox}{Lessons Learned From the Unit-Field Walk}
The component should never pretend that formatted text and stored value are the
same fact. Keep draft text, parsed value, selected unit, canonical value, and
problems distinct. That gives the UI room to be humane while giving the
application a precise value at commit time.
\end{principlebox}

\section{Advanced JTable Column Headers}

\begin{roadmapbox}
This section extends table headers from labels into stateful controls. It covers
sort glyphs, popup menus, visibility commands, filter entry points, help,
descriptor-backed identity, and persistence. The recurring warning is that
localized header text is display data, not a reliable key for behavior.
\end{roadmapbox}

Table headers in business applications do more than display text. They show sort
direction, open column menus, route help, explain validation, expose filtering entry
points, and manage column visibility. Existing table chapters cover model indexes,
descriptors, selection, table state, rendering, and large data. This section focuses on
the header as a practical compound surface. The header is part renderer, part command
surface, part metadata router, and part persisted-state controller.

The first rule is unchanged from the table chapters: do not derive identity from visible
header text. A header cell represents a descriptor column, not a localized string. The
current view column can move. The current model column can differ under reordering. The
visible text can change with locale. The header behavior should map a mouse or keyboard
location to a \api{ColumnKey} or descriptor, then use that stable fact for sorting,
menus, help, and state.

\begin{figure}[htbp]
\centering
\includegraphics[width=0.96\linewidth]{\practicalTableHeaderFigure}
\caption{An advanced table header is a descriptor-backed command surface. Pointer location resolves to column identity before sort, filter, help, or visibility commands run.}
\label{fig:practical-table-header}
\end{figure}

\begin{lstlisting}[caption={Header commands begin by resolving a descriptor column.}]
ColumnDescriptor<CustomerRow, ?> columnAt(Point point) {
    JTableHeader header = table.getTableHeader();
    int viewColumn = header.columnAtPoint(point);
    if (viewColumn < 0) {
        return null;
    }

    int modelColumn = table.convertColumnIndexToModel(viewColumn);
    return tableModel.descriptor().columnAt(modelColumn);
}
\end{lstlisting}

The example is short but important. Header behavior starts with coordinates and ends
with descriptor metadata. Everything else should use the descriptor. If a popup item
hides a column, it hides by persistence id. If F1 opens help, it opens the descriptor
help topic. If a tooltip is shown, it resolves the resource key from the descriptor. If
a filter row is opened, it uses the column value type and filter capability. The header
string is only display.

Sorting has two parts: user gesture and visual state. The gesture may be a click,
keyboard command, menu item, or programmatic action. The visual state is the current
sort key list in
\api{RowSorter}. The header renderer
should paint sort glyphs from sorter state, not remember its own last click. Multi
column sorting, toggling direction, clearing sort, and restoring persisted sort state
all become easier when the renderer observes the sorter rather than owning sorting.

\begin{lstlisting}[caption={A header renderer reads sort state rather than storing it.}]
TableCellRenderer sortingHeaderRenderer(
        TableCellRenderer delegate,
        JTable table) {
    return (tbl, value, selected, focus, row, viewColumn) -> {
        Component c = delegate.getTableCellRendererComponent(
                tbl, value, selected, focus, row, viewColumn);

        int modelColumn = table.convertColumnIndexToModel(viewColumn);
        SortOrder order = sortOrderFor(table.getRowSorter(), modelColumn);
        decorateHeader(c, order);
        return c;
    };
}
\end{lstlisting}

The renderer can decorate with an icon, text suffix, or custom label. It should remain
cheap and side-effect free. It should not change sorter state, open popups, or install
listeners. Rendering is observation. Mouse and keyboard commands perform mutation.
Keeping that separation prevents a surprising class of bugs where repainting a header
changes table state.

Context menus should be constructed from column capabilities. A numeric column may offer
range filter. A date column may offer date filter and quick periods. A text column may
offer contains, starts with, and clear filter. A hideable column may offer Hide Column.
A required column may disable hiding and explain why. A column with help metadata may
offer Help. The header menu should not be one global menu with every item enabled for
every column.

\begin{lstlisting}[caption={A header popup receives a column descriptor.}]
JPopupMenu createHeaderMenu(ColumnDescriptor<?, ?> column) {
    JPopupMenu menu = new JPopupMenu();
    menu.add(new JMenuItem(sortAscendingAction(column.key())));
    menu.add(new JMenuItem(sortDescendingAction(column.key())));
    menu.add(new JMenuItem(clearSortAction(column.key())));
    menu.addSeparator();

    if (column.isHideable()) {
        menu.add(new JMenuItem(hideColumnAction(column.key())));
    }
    if (column.helpTopic().isPresent()) {
        menu.add(new JMenuItem(helpAction(column.helpTopic().get())));
    }
    return menu;
}
\end{lstlisting}

The actual descriptor API may differ. The structure should not. The popup is created or
refreshed with stable metadata. Actions receive keys or topics. The menu item text comes
from resources. The menu does not parse the visible header label and does not store view
indexes for later use. If the user opens a popup and then a column moves before a
command runs, the command should still refer to the original descriptor or re-resolve
deliberately.

Popup triggers need careful selection policy. A table body popup often selects the row
under the pointer before showing row commands. A header popup usually should not alter
row selection. It operates on columns. If the table supports focused column state or
header keyboard focus, opening the menu may update the focused column. That state should
be separate from row selection. Users should not lose their current row selection
because they opened a column visibility menu.

Keyboard access for headers is often neglected. Standard \api{JTable} header focus
behavior is limited. Serious applications can add commands to the table or surrounding
view: sort focused column, open header menu for focused column, clear sort, show filter
panel, and show help. The focused column may come from selected cell, lead column, or a
separate header focus model. Pick one and document it. Do not require a mouse for column
menus that control core table behavior.

Column visibility belongs to table state. Hiding a column should remove or hide the
\api{TableColumn} in the
\api{TableColumnModel}; it should not
remove the column from the underlying table model unless the data schema really changes.
Restoring visibility should use the descriptor and persisted order. The header menu is
merely a command surface for table-state policy. This keeps export, help, validation,
and model identity coherent.

\begin{pitfallbox}{Header Text Persistence}
Persisting widths, order, visibility, filters, or sort by visible header text will break
under localization and copy editing. Persist by descriptor id or column key.
\end{pitfallbox}

Filter entry points should avoid turning the header into a second table model. A filter
icon in the header can open a filter editor, show a filter summary, or clear the current
filter. The filter model itself belongs to the table presenter or search state. The
header renderer can show whether a filter is active. The popup can modify filter state
through commands. It should not store filter text in renderer components.

Tooltips and help use the same metadata. Header tooltip binding can resolve short
explanatory text from resources. F1 on a header or menu item can route to a help topic.
A filtered column may add a tooltip line that summarizes the current filter. A sorted
column may mention sort direction. These are presentation facts derived from descriptor
and table state. They should be composed in a policy object rather than hard-coded in a
renderer.

Accessibility should make sort and menu state audible. A header for a sorted column
should expose the column name and sort direction. A header with a menu should expose
that a popup is available. A hidden-column command should have names that match
localized column labels but should operate on stable ids. If the header uses icons for
sort or filter state, the same state must be available as accessible description or
text.

Review an advanced header by following a right-click. Pointer location maps to view
column, then model column, then descriptor. Descriptor supplies column id, resource
keys, help, hideability, value type, and persistence id. The menu builds actions from
those facts. Actions mutate sorter, filter, or table state. Renderer observes the
resulting state. Persistence writes stable ids. Cleanup removes installed listeners when
the view closes. If any step switches back to header text, the design is incomplete.

\begin{principlebox}{Lessons Learned From Advanced Headers}
A useful header is a table control surface with stable column identity. Sorting,
filtering, help, visibility, and persistence should all route through descriptor
keys and table state. The renderer paints current state; it should not become
the place where column behavior is discovered.
\end{principlebox}

\section{Design walk: header state, descriptor state, and Swing state}

\begin{roadmapbox}
This walk separates three states that often get mixed together: Swing's view
state, the table model's domain state, and Corusco descriptor state. It traces a
header click and a header popup command through those layers so the same column
keeps its identity after sorting, reordering, hiding, or localization.
\end{roadmapbox}

An advanced table header crosses three coordinate systems. The pointer or keyboard focus
is in header view coordinates. The \api{JTable} maps view columns to model columns.
Corusco table descriptors map model columns to stable column keys and metadata. The
header component should cross those systems deliberately every time. Storing a view
column inside a menu action is only safe if the action runs immediately and the column
cannot move before it runs. Storing a column key is safer and easier to test.

The header model should not duplicate table state. Sort state belongs to the
\api{RowSorter} or presenter sort model. Visibility belongs to the table-state
controller. Filter state belongs to a filter model. Help belongs to descriptor metadata.
Tooltip state is composed from descriptor, sort, filter, and problems. The header
renderer and popup observe these facts. They should not maintain private copies that can
diverge.

\begin{lstlisting}[caption={Header context names the resolved table facts.}]
record HeaderContext(
        TableKey tableKey,
        ColumnKey columnKey,
        int viewColumn,
        int modelColumn,
        ColumnDescriptor<?, ?> descriptor) {
}
\end{lstlisting}

A context object like this is useful even if it exists only inside the header installer.
It can be logged, tested, and passed to action factories. A Sort Ascending action can
accept
\api{columnKey}. A Help action can
accept the descriptor's help topic. A Hide action can accept the column persistence id.
The context makes coordinate conversion visible and prevents late code from
rediscovering it incorrectly.

The event path for a header popup is precise. Mouse event arrives on the EDT. Installer
checks whether it is a popup trigger. It resolves
\api{HeaderContext}
from point. It asks the presenter or table command factory for menu actions. It
refreshes menu item state from current sort, filter, and visibility. It shows the popup
at the header. Choosing a menu item invokes an action that mutates sorter, filter model,
table state, or help service. The renderer later observes changed state and repaints. At
no point should the popup mutate row selection unless the product has a separate
focused-column model and names it.

\begin{lstlisting}[caption={A header installer converts events before action creation.}]
void maybeShowHeaderPopup(MouseEvent event) {
    if (!event.isPopupTrigger()) {
        return;
    }
    HeaderContext context = resolveHeaderContext(event.getPoint());
    if (context == null) {
        return;
    }
    JPopupMenu menu = headerMenuFactory.create(context);
    menu.show(header, event.getX(), event.getY());
}
\end{lstlisting}

Corusco makes the menu factory more useful because descriptor columns already carry
stable ids, resource keys, persistence ids, and capability flags. A plain Swing table
can still implement the pattern, but Corusco removes the need to build a parallel
metadata map keyed by integer column indexes. The descriptor is the map. The header
installer can be generic over row type because it does not know business meaning; it
only passes descriptor facts to actions.

The renderer needs a different kind of discipline. It receives a value and view column.
It may ask for sort direction, filter active state, and maybe problem indicators. It
should render quickly and side-effect free. It should not build popup menus, install
listeners, or persist state. In a large table, header renderers are cheap compared with
cell renderers, but the same rule applies: rendering is presentation of existing state,
not computation of new state.

Failure modes are mostly coordinate bugs. A user right-clicks Due Date, then columns
move and Hide hides Amount. A tooltip says the localized header text is the help topic.
Persisted visibility breaks after translation. A filter indicator remains painted after
filter state clears because renderer cached it privately. A popup action reads selected
row and accidentally operates on body selection. The cure is the same every time:
resolve context once, pass stable ids, let owners mutate their own state, and make
renderers observe.

\begin{principlebox}{Lessons Learned From the Header-State Walk}
The table header must translate from view coordinates to descriptor identity
before it makes decisions. That translation is the difference between a header
that survives reordering and localization and one that only works in the initial
demo screen.
\end{principlebox}

\section{Temporal Selection Panels}

\begin{roadmapbox}
This section designs calendar, date-time, and date-time span selectors. It
starts with value types and zone policy, then moves through draft versus
committed values, disabled days, min/max bounds, quick ranges, keyboard
navigation, popup commit rules, and the validation problems that appear at time
boundaries.
\end{roadmapbox}

Temporal selection controls look simple because everyone has used a calendar. They are
difficult because time carries policy. A date may be local to the business. A date-time
may be local to a branch, user, or server. An instant may be stored in UTC but displayed
in a zone. A range may be inclusive at the start and exclusive at the end. A date may be
nullable, constrained, disabled, or only valid on business days. A calendar control
should not hide these policies behind painted day cells.

Use Java time types deliberately.
\api{LocalDate} is a calendar date
without time or zone. \api{LocalDateTime} is a date and time without zone; it is not an
instant. \api{Instant} is a point on the timeline. \api{ZoneId} is the policy that
converts between local display and instant storage. A date picker for invoice due date
may commit
\api{LocalDate}. A meeting scheduler
may commit \api{ZonedDateTime} or an
\api{Instant} plus zone. A log filter
may commit an instant range. The component must know which contract it serves.

\begin{figure}[htbp]
\centering
\includegraphics[width=0.96\linewidth]{\practicalTemporalFigure}
\caption{Temporal selectors need visible policy. Draft date, draft time, quick ranges, zone, endpoint semantics, Apply, and Cancel are separate facts even when presented in one panel.}
\label{fig:practical-temporal-selection}
\end{figure}

\begin{tabularx}{0.96\linewidth}{@{}p{0.24\linewidth}X@{}}
\toprule
Value type & Use in business UI\\
\midrule
\api{LocalDate} & Due dates, service dates, birthdays, accounting periods, and other date-only facts.\\
\api{LocalDateTime} & Local appointment-like entries when zone is supplied by surrounding context.\\
\api{Instant} & Audit timestamps, server events, import windows, and stored timeline facts.\\
\api{DateRange} & Inclusive or exclusive date intervals with named endpoint policy.\\
\api{DateTimeSpan} & Start and end date-time values plus zone and endpoint policy.\\
\bottomrule
\end{tabularx}

A normal calendar picker has at least four states: visible month, focused day, draft
selected day, and committed value. Visible month changes when the user pages. Focused
day changes with keyboard navigation. Draft selected day may change when the user clicks
or arrows. Committed value changes when the user accepts. Some small date pickers commit
immediately on click; larger panels often separate draft and commit so Cancel can
restore the old value. The choice should match the surrounding workflow.

\begin{lstlisting}[caption={Calendar state separates visible month, focus, and value.}]
record CalendarState(
        YearMonth visibleMonth,
        LocalDate focusedDate,
        LocalDate draftSelection,
        LocalDate committedValue) {

    CalendarState focus(LocalDate date) {
        return new CalendarState(
                YearMonth.from(date), date, draftSelection, committedValue);
    }
}
\end{lstlisting}

The sketch is immutable for clarity. A real component may use observable values, mutable
fields confined to the EDT, or a presenter model. The separation matters. Paging from
March to April should not commit a value. Moving focus with arrow keys should not
necessarily commit a value. Pressing Enter may commit the focused date. Pressing Escape
may restore the committed value and close the popup. These behaviors are impossible to
review if the component has only a single
\api{selectedDate} field.

Calendar grids are good candidates for custom painting or renderer-based panels. A
simple implementation can use a \api{JPanel} with seven columns and day buttons. That
gives keyboard and accessibility help but creates many components. A custom
\api{JComponent} can paint the grid and handle hit testing, which is efficient and
visually flexible, but it must implement focus, keyboard, accessible names, tooltips,
and repaint bounds carefully. Choose based on product needs. A dense business calendar
with disabled days, badges, and range previews may justify custom painting. A small
picker may not.

Disabled days should be a policy, not hard-coded painting. The calendar can ask a
predicate whether a date is selectable and another service for day metadata such as
holiday, business closure, warning, or event count. The predicate should be cheap and
deterministic for visible months. If disabled state requires server data, load it before
or around the picker, show loading state, and keep the component itself responsive. A
date cell renderer should not call a remote service during paint.

\begin{lstlisting}[caption={Selectable-date policy stays outside painting.}]
interface DateSelectionPolicy {
    boolean isSelectable(LocalDate date);
    Optional<String> disabledReason(LocalDate date);
}
\end{lstlisting}

Keyboard navigation is part of the calendar contract. Left and right move by one day. Up
and down move by one week. Page Up and Page Down move by month. Home and End may move to
week boundaries. Ctrl+Home may move to today or the start of the visible month depending
on product convention. Enter commits. Escape cancels or closes. The component should
document and test the chosen bindings. These are not extras; they are how many users
operate date fields quickly.

A date-time picker combines a date selection with time entry. Do not hide the time
policy inside a formatted text field. Time may have minute increments, business hours,
disabled intervals, seconds, timezone display, and default time when a date is chosen. A
practical design uses a calendar for date and a separate time editor or time list for
time. The committed value is produced only after date, time, and zone policy agree.

\begin{lstlisting}[caption={Date-time draft state names the zone policy.}]
record DateTimeDraft(
        LocalDate date,
        LocalTime time,
        ZoneId displayZone,
        Optional<Instant> committedInstant,
        List<Problem> problems) {
}
\end{lstlisting}

If the application stores
\api{Instant}, the component should
convert through the explicit \api{ZoneId}. If the application stores
\api{LocalDateTime}, the
component should still know where the zone policy lives, even if it is outside the
committed value. Daylight-saving transitions can make local times invalid or ambiguous.
Business applications often choose a simple policy: reject invalid local times and
choose the earlier or later offset for ambiguous times. Whatever the policy is, document
it. Do not leave it to whatever
\api{atZone} happens to do in one code path.

Date-time span selectors add another layer. The user chooses start and end of an
interval. The UI may show two calendars, one calendar plus fields, quick range presets,
duration fields, and validation summary. The model should name start draft, end draft,
zone, endpoint policy, and quick preset. It should not store only two strings. Problems
such as "end before start", "range too long", "start outside business hours", and "end
required" need structured targets.

\begin{lstlisting}[caption={A date-time span has endpoint policy as data.}]
record DateTimeSpanDraft(
        LocalDateTime start,
        LocalDateTime end,
        ZoneId zone,
        EndpointPolicy endpointPolicy,
        Optional<QuickRange> quickRange,
        List<Problem> problems) {

    boolean ordered() {
        return !end.isBefore(start);
    }
}
\end{lstlisting}

Endpoint language matters. "From June 1 to June 30" may mean inclusive dates. An instant
interval for database queries often means start inclusive, end exclusive. A date-time
span for shift scheduling may include both endpoints visually but store an exclusive end
instant. Users should not have to infer this from code. Labels, help text, and
validation should use the same policy as storage. Tests should cover boundary values
because off-by-one temporal bugs are easy to ship.

Quick ranges are commands, not magic values. "Today", "This week", "Last 30 days",
"Current quarter", and "Custom" all set or describe span state. If a user edits start or
end manually after choosing "Last 30 days", the quick range may become Custom. If the
current date changes at midnight while the screen is open, a relative range may need
recomputation or a clear capture time. A quick range command should state whether it is
evaluated when chosen, when submitted, or continuously. Most business screens evaluate
when chosen and let the user review the resulting dates.

Draft and commit boundaries are essential for temporal popups. A small date field popup
may commit immediately and close on day click. A date-time span panel with quick ranges
and validation should usually have Apply and Cancel. Cancel restores the committed span.
Apply validates the draft and commits if valid. Closing the popup by Escape should
behave like Cancel, not like Apply. Clicking outside the popup needs a product decision.
The implementation should make that decision explicit rather than relying on default
popup dismissal to lose state silently.

Accessibility for calendar controls is work, but it is manageable if the model is clear.
Each day should have a name such as "Tuesday, June 16, 2026", state such as selected,
focused, disabled, today, or outside current month, and a disabled reason if applicable.
A span preview should not be color-only. Time fields should expose format and valid
increments. Quick range buttons should say what dates they select. A custom-painted
calendar needs extra care because Swing will not automatically expose painted day cells
as accessible children.

Temporal controls should be tested with fixed clocks. A test that depends on today's
actual date will fail later for no useful reason. Inject
\api{Clock}
or pass the current date into the model. Test month boundaries, leap days,
daylight-saving transitions if the control handles instants, disabled days, range
ordering, cancel behavior, and keyboard navigation. Screenshot tests should include a
month that starts late in the week and a month with six visible rows so grid geometry is
exercised.

\begin{principlebox}{Lessons Learned From Temporal Selectors}
Temporal controls are mostly policy controls. Their visible calendar grid is
only the surface over value type, zone, bounds, disabled dates, endpoint
semantics, draft state, and commit behavior. Treat those policies as model data
so date-time bugs are not hidden in mouse listeners.
\end{principlebox}

\section{Design walk: temporal models and commit boundaries}

\begin{roadmapbox}
This walk focuses on the moment where a temporal picker changes from draft to
committed value. It compares date, date-time, instant, and interval models,
shows how quick ranges and endpoint policies affect storage, and explains how
cancel, apply, and focus loss should leave the model in predictable states.
\end{roadmapbox}

Temporal selectors force the application to choose between immediate commit and draft
commit. A small date popup in a simple form may commit immediately when the user clicks
a day. A date-time span panel should almost never commit on every click because the user
may need to adjust start, end, time, quick range, and zone before the value is
meaningful. A dialog may have OK and Cancel outside the selector; an inline report
filter may have Apply inside the selector. The model should name the current draft and
the last committed value so Cancel and dirty state are unambiguous.

\begin{lstlisting}[caption={Temporal editor state separates draft from committed value.}]
record TemporalEditorState<T>(
        T committedValue,
        T draftValue,
        boolean dirty,
        ProblemSet problems) {

    TemporalEditorState<T> cancel() {
        return new TemporalEditorState<>(
                committedValue, committedValue, false, ProblemSet.empty());
    }
}
\end{lstlisting}

The generic type might be
\api{LocalDate}, \api{DateTimeDraft},
or
\api{DateTimeSpanDraft}. The exact type matters less than the separation.
Corusco field models already use this style: raw text and parsed value are not the same.
Temporal selectors need the same discipline because navigating a calendar, changing a
time spinner, and selecting a quick range are draft edits until the workflow says
otherwise.

The next decision is clock ownership. A calendar often highlights today and quick ranges
often depend on today. If the component calls
\api{LocalDate.now()} directly, tests become time-dependent and long-lived
screens can become stale at midnight. Prefer a \api{Clock} or a presenter supplied
"today" value. Quick ranges can then be evaluated at a known instant. The product can
decide whether an open screen updates at midnight or keeps the captured date until the
user refreshes.

Event wiring for a calendar is richer than it looks. Arrow key changes focus. Page key
changes visible month and focus. Mouse hover may change preview state but not selection.
Mouse click changes draft selection. Double click or Enter may commit. Escape cancels.
Month buttons change visible month without changing committed value. Disabled dates
block selection and may show a reason. These events should call model methods with names
like
\api{focusDate}, \api{pageMonth}, \api{selectDraft}, \api{commitDraft}, and
\api{cancelDraft}. A mouse listener that directly sets \api{selectedDate}
cannot express these differences.

\begin{lstlisting}[caption={Calendar events should call named model operations.}]
void handleDayClick(LocalDate date) {
    if (!policy.isSelectable(date)) {
        model.focusDate(date);
        model.showDisabledReason(date);
        return;
    }
    model.selectDraft(date);
    if (commitOnClick) {
        model.commitDraft();
    }
}
\end{lstlisting}

Corusco helps by giving temporal selectors the same binding shape as other fields. A
generated \api{@DateField} can cover stock date fields. A custom date-time span selector
can still expose a binding interface that writes to a field model, publishes problems,
and accepts resources and help topics. The component can be custom; the surrounding form
still sees a field key, problem target, dirty state, and commit result. That is the
maintainability point.

Date-time selectors need a zone decision at model construction time. If the zone comes
from user preferences, it should be part of the field context. If the zone comes from
the selected branch or account, a change in that context may invalidate the draft. If
the storage type is \api{Instant}, conversion errors around daylight-saving transitions
should produce problems, not silent adjustments. If the storage type is
\api{LocalDateTime}, the model should
still document what zone will be used later when the value becomes an instant.

Span selectors also need endpoint semantics in both model and text. A report query with
start inclusive and end exclusive should not label the end as "through" if users
interpret that as inclusive. A date-only UI may display "June 1 through June 30" but
convert to instant interval
\api{[June 1 00:00, July 1 00:00)}. That conversion should be in one place,
tested with boundary dates. A useful model stores \api{EndpointPolicy} so help,
validation, and storage agree.

Failure modes are often expensive. A report misses the last day because UI and database
disagree about exclusive end. A daylight-saving gap is silently shifted by an hour. A
quick range named "Last 7 days" changes while the user is reviewing it because the clock
changed. Cancel closes a popup but leaves draft values in the field model. A disabled
holiday can be clicked by keyboard even though mouse selection blocks it. These are not
calendar painting bugs. They are model and event-path bugs, so the model and event path
need first-class design.

\begin{principlebox}{Lessons Learned From Temporal Model Boundaries}
Commit and cancel rules are as important as calendar navigation. Users can
tolerate a picker that pages months differently from another product, but they
cannot tolerate a picker that silently changes endpoint semantics, loses a draft
on focus transfer, or stores an instant in an undocumented zone.
\end{principlebox}

\section{Tag Cloud Component}

\begin{roadmapbox}
This section builds a tag cloud as a business editor, not as a decorative word
cloud. The path is: define the value and query model, choose whether tags are
free-form or catalog-backed, decide how wrapping and focus work, wire keyboard
and mouse events through commands, add asynchronous suggestions without races,
bind validation through Corusco problem targets, and finish with review rules
for dense screens.
\end{roadmapbox}

Tag clouds appear in business applications under many names: labels, categories,
recipients, skills, product attributes, risk flags, markets, recipients, and
saved filters. The visual form is familiar. Selected tags appear as compact
chips, each chip has a label and usually a remove affordance, and a small editor
at the end lets the user add more tags. The component may show suggestions while
typing, may color tags by category, may allow drag reordering, and may limit how
many tags can be selected. It looks like a small convenience widget, but it is a
surprisingly rich editor.

The first design decision is whether the cloud edits strings, catalog entries,
or typed domain references. A string tag such as \texttt{"urgent"} is easy to
display but hard to merge, localize, validate, or audit. A catalog tag has a
stable id, display label, color policy, and perhaps a disabled reason. A domain
reference tag may point to a user, cost center, product, or jurisdiction. The
component can display all of them as chips, but the model should not collapse
them into display text. A label is for the user. Identity is for the application.

\begin{figure}[htbp]
\centering
\includegraphics[width=0.96\linewidth]{\practicalTagCloudFigure}
\caption{A tag cloud is a list editor with a text-entry tail. Stable tag ids,
suggestion state, focused chip, category styling, and validation problems all
belong to the model even when the visible control is built from ordinary Swing
children.}
\label{fig:practical-tag-cloud}
\end{figure}

\begin{principlebox}{Tag Identity Is Not Tag Text}
Never use the rendered chip label as the selected value. Use a stable tag id or
domain key, then derive label, color, tooltip, and help from resources or
descriptor metadata.
\end{principlebox}

A practical tag model has at least five pieces of state. It has the committed or
draft list of selected tags. It has the query currently typed by the user. It
has a suggestion list, including loading and failure state if suggestions come
from a service. It has focus state, because keyboard users need to move between
chips and the editor. It has problems, because duplicates, incompatible tags,
too many tags, and disabled tags are validation issues rather than painting
issues. More advanced models also hold a maximum count, ordering policy, color
policy, and source scope.

\begin{lstlisting}[caption={A tag cloud model separates identity, query, suggestions, and focus.}]
record TagCloudDraft(
        List<TagSelection> selected,
        String query,
        SuggestionState<TagSuggestion> suggestions,
        FocusTarget focus,
        List<Problem> problems) {

    boolean canAdd(TagSuggestion suggestion) {
        return problems.stream().noneMatch(Problem::blocksCommit)
                && selected.stream().noneMatch(tag -> tag.sameId(suggestion));
    }
}

record TagSelection(
        TagId id,
        String label,
        TagCategory category,
        boolean removable) {
}

sealed interface FocusTarget {
    record Editor() implements FocusTarget {}
    record Chip(TagId id) implements FocusTarget {}
}
\end{lstlisting}

The explicit focus target is worth the ceremony. Without it, the component often
tries to infer focus from the child component that happens to own keyboard focus
or from the last painted chip rectangle. Both choices fail when the chip list
wraps, when a chip is removed, or when a screen reader needs a stable target. If
focus belongs to a model value, removing the focused chip can move focus to a
neighbor by rule. Repainting then follows the rule instead of inventing a rule.

The next decision is the component shape. A small tag cloud can be a
\api{JPanel} with a wrapping layout, chip buttons, and a \api{JTextField}.
This is often the best first implementation. It inherits focus painting, button
semantics, tooltips, font metrics, and accessibility from Swing. The panel
should not use \api{FlowLayout} blindly, because a text editor at the end needs
a useful minimum width and a predictable baseline. A custom wrapping layout or a
MigLayout with wrap constraints usually produces more stable geometry.

\begin{lstlisting}[caption={A tag cloud can start as a panel of chip buttons plus an editor.}]
final class TagCloudPanel extends JPanel {
    private final TagCloudModel model;
    private final JTextField editor = new JTextField();
    private final Map<TagId, JComponent> chips = new LinkedHashMap<>();

    TagCloudPanel(TagCloudModel model) {
        super(new WrapRowLayout(6, 6));
        this.model = model;
        setBorder(UIManager.getBorder("TextField.border"));
        setBackground(UIManager.getColor("TextField.background"));
        installEditorActions();
        model.changes().subscribe(this::refresh);
        refresh(model.snapshot());
    }

    private void refresh(TagCloudDraft draft) {
        removeAll();
        chips.clear();
        for (TagSelection tag : draft.selected()) {
            JComponent chip = createChip(tag, draft.focus());
            chips.put(tag.id(), chip);
            add(chip);
        }
        editor.setText(draft.query());
        add(editor);
        revalidate();
        repaint();
    }
}
\end{lstlisting}

This sketch is deliberately ordinary Swing. The important part is the refresh
boundary. The view is a projection of the model. Child chip components may be
rebuilt when the model changes because their identity is the tag id. For a very
large cloud, a diffing refresh is better. Most business tag fields contain fewer
than thirty selected tags, so a simple rebuild is easier to review and perfectly
fast. The model remains the owner of selected tags and focus, so rebuilding does
not lose important state.

Chip components should be command surfaces rather than passive labels with mouse
listeners. A chip has a focus action, a remove action, perhaps a details action,
and maybe a context menu. A remove icon should invoke the same command as Delete
when the chip is focused. If a tag is not removable, the remove command should be
disabled with a reason that can be shown as tooltip text. Corusco command
metadata is useful here because a toolbar button, popup menu item, and chip
affordance can all share one command contract.

\begin{lstlisting}[caption={A chip remove action carries the tag id and disabled reason.}]
final Action removeTag(TagSelection tag) {
    Command command = commands.removeTag(tag.id());
    return new AbstractAction(resources.text(command.labelKey())) {
        {
            putValue(SHORT_DESCRIPTION, command.description(resources));
            setEnabled(command.enabled());
        }

        @Override
        public void actionPerformed(ActionEvent event) {
            model.remove(tag.id());
        }
    };
}
\end{lstlisting}

Keyboard behavior is the design center of a good tag cloud. The editor accepts
normal text editing. Enter adds the highlighted suggestion or converts the query
into a free-form tag if policy permits. Escape closes suggestions or clears the
query depending on state. Backspace in a non-empty query remains a text-editing
operation. Backspace in an empty query moves focus to the previous chip or
removes the focused chip on a second press, depending on the product convention.
Left and right move between chips and the editor in component orientation order.
Delete removes the focused chip. Ctrl+Space or Alt+Down opens suggestions.

\begin{tabularx}{0.96\linewidth}{@{}p{0.25\linewidth}X@{}}
\toprule
Gesture & Model operation\\
\midrule
Enter in query & Accept highlighted suggestion or create free-form tag.\\
Escape & Close suggestions, then clear query, then leave field by normal focus rule.\\
Backspace in empty query & Focus previous chip or remove already focused chip.\\
Left/Right & Move focus by visual order, respecting component orientation.\\
Delete on focused chip & Invoke remove command with disabled reason if blocked.\\
Alt+Down & Open suggestion popup from the editor baseline.\\
\bottomrule
\end{tabularx}

Suggestions add another design decision: synchronous catalog filtering or
asynchronous search. If the tag catalog is small and already local, the model can
filter in memory. If suggestions come from a service, every query change starts
a request that may return after a later request. The component must ignore stale
results. A simple generation number is enough. Corusco lifecycle scopes help
because the subscription and background task can be cancelled when the screen
closes. Without cancellation, a stale suggestion callback can update a disposed
view or reopen a popup after the form is gone.

\begin{lstlisting}[caption={Suggestion loading uses a generation token to ignore stale replies.}]
final class TagSuggestionController {
    private final AtomicLong generation = new AtomicLong();
    private final TagCloudModel model;
    private final TagService service;

    void queryChanged(String query) {
        long request = generation.incrementAndGet();
        model.setSuggestions(SuggestionState.loading(query));
        service.searchTags(query).whenComplete((items, failure) -> {
            if (request != generation.get()) {
                return;
            }
            if (failure != null) {
                model.setSuggestions(SuggestionState.failed(query, failure));
            } else {
                model.setSuggestions(SuggestionState.ready(query, items));
            }
        });
    }
}
\end{lstlisting}

Free-form tags require a policy object, not a Boolean scattered through key
handlers. Some screens permit any value. Some permit a new tag only when the user
has a permission. Some require a prefix, minimum length, or catalog review. Some
convert case or trim punctuation. The model should expose the policy result as a
problem or command state. If the user presses Enter and the query cannot become a
tag, the component should keep focus, select the problematic part if useful, and
expose the problem text through the same validation system as the rest of the
form.

\begin{lstlisting}[caption={Creation policy returns a domain result, not a UI decision.}]
sealed interface TagCreationResult {
    record Allowed(TagSelection tag) implements TagCreationResult {}
    record Rejected(Problem problem) implements TagCreationResult {}
    record NeedsConfirmation(TagCandidate candidate) implements TagCreationResult {}
}
\end{lstlisting}

Layout is less trivial than it looks. A tag cloud usually wraps, but wrapping
changes the field height. In a dense form, that height change can move unrelated
controls while the user types. There are three practical policies. The control
can grow vertically, which is best for editing but can make forms jump. It can
have a maximum row count and scroll internally, which keeps the form stable but
adds a nested scrolling surface. It can collapse chips into a summary after a
limit, which is useful for read-mostly filters. The policy should be visible in
the component contract because it changes validation placement and screenshot
expectations.

\begin{detailbox}{Baseline and Wrapping}
Swing baseline alignment is easy for a single \api{JTextField} and much harder
for a wrapping composite. If a tag cloud appears in a form row, decide whether
the row aligns to the first chip baseline, the editor baseline, or the top edge.
Then encode that in the layout manager or wrapper panel instead of relying on
the current font.
\end{detailbox}

Painting should communicate category without making color the only signal. A
priority tag may be red, but it also needs text, icon, or tooltip semantics. A
disabled tag should look disabled and explain why it cannot be removed or
selected. A selected or focused chip needs a focus ring that passes high-contrast
themes. Use UI defaults for focus colors and borders where possible. If the
application uses FlatLaf client properties for style variants, keep those
properties in a chip factory rather than repeating them in every screen.

Corusco helps the tag cloud in three places. First, generated field keys provide
stable problem targets: the list can have a problem, and a specific tag id can
have a problem. Second, generated resources and descriptors keep tag labels,
tooltips, disabled reasons, and help topics out of the component. Third,
lifecycle scopes own suggestion subscriptions and popup listeners. The component
itself can remain a small Swing view with a model interface. The Corusco binding
connects that view to generated form state.

\begin{lstlisting}[caption={A Corusco binding maps tag model problems to the cloud view.}]
final class TagCloudBinding implements Binding {
    private final FieldKey<List<TagId>> fieldKey;
    private final TagCloudPanel view;
    private final TagFieldModel field;
    private final ProblemBinder problems;
    private final Subscription subscription;

    TagCloudBinding(
            FieldKey<List<TagId>> fieldKey,
            TagCloudPanel view,
            TagFieldModel field,
            ProblemBinder problems) {
        this.fieldKey = fieldKey;
        this.view = view;
        this.field = field;
        this.problems = problems;
        this.subscription = field.changes().subscribe(this::refresh);
        view.addTagCloudListener(this::viewChanged);
        refresh(field.snapshot());
    }

    private void viewChanged(TagCloudEvent event) {
        field.setDraft(event.selectedIds());
    }

    private void refresh(TagFieldSnapshot snapshot) {
        view.setDraft(snapshot.toCloudDraft());
        problems.show(fieldKey, snapshot.problems(), view);
    }

    @Override
    public void close() {
        subscription.close();
    }
}
\end{lstlisting}

The event object should describe intent, not pixels. Good events are
\texttt{queryChanged}, \texttt{tagAccepted}, \texttt{tagRemoved},
\texttt{chipFocused}, and \texttt{suggestionRequested}. Poor events are
\texttt{mouseClickedAt}, \texttt{labelRemoved}, or \texttt{popupClosed}
when those are the only facts the presenter receives. The view may internally
handle a mouse click, but the model and binding should see the domain operation.
This keeps unit tests short and prevents a screen from depending on chip
geometry.

Validation is often richer than "required." A tag may be unknown, duplicated,
blocked by permission, incompatible with another tag, expired in the catalog, or
valid only in a chosen scope. The model should decide whether invalid tags stay
visible. In many business screens they should, because silently dropping a tag
destroys user context. A blocked tag can remain as a chip with a problem marker
and disabled remove state if the screen must preserve historical data. A new
entry form may instead reject the tag immediately. Both policies are valid, but
the policy belongs in the model and review checklist.

Popup suggestions need careful ownership. The popup should be positioned from
the editor or from the whole field, depending on whether suggestions relate to
the query or to the selected set. It should close on commit, Escape, focus loss
outside the cloud, model disposal, and Look-and-Feel update. It should not close
when focus moves from the editor into the suggestion list if that list is part of
the popup interaction. These details are much easier to reason about when popup
state is represented as part of the suggestion model.

\begin{tabularx}{0.96\linewidth}{@{}p{0.23\linewidth}X@{}}
\toprule
Review area & Acceptance rule\\
\midrule
Identity & Selected values are stable ids or domain references, never display text.\\
Keyboard & Every mouse operation has a named keyboard operation.\\
Async & Stale suggestion results cannot replace newer query results.\\
Validation & Duplicate, disabled, incompatible, and over-limit tags have problem targets.\\
Layout & Wrapping policy is explicit and tested at large font sizes.\\
Lifecycle & Suggestion tasks, popup listeners, and model subscriptions close with the screen.\\
\bottomrule
\end{tabularx}

\begin{principlebox}{Lessons Learned From Tag Clouds}
A tag cloud is a list editor with a text editor attached. Its hard problems are
identity, focus, suggestion lifetime, and validation, not chip painting. Corusco
is most valuable when it supplies stable field keys, resource-backed labels,
problem targets, command metadata, and lifecycle cleanup so the Swing component
can remain a projection of a reviewable model.
\end{principlebox}

\section{Interactive Cells in JTable}

\begin{roadmapbox}
This section designs cells that look and behave interactive inside a
\api{JTable}. The path is: separate rendering from editing, choose row identity,
route commands through descriptors, convert view and model coordinates safely,
decide when a real editor is needed, handle keyboard and focus, coordinate with
sorting and filtering, and use Corusco table metadata to avoid deriving behavior
from localized cell text.
\end{roadmapbox}

Interactive table cells are common in business screens. A row has a status pill
that opens a history popup. A cell contains a small "Approve" action. A quantity
cell has plus and minus buttons. A progress cell exposes retry. A boolean cell is
a tri-state editor. A row action column contains icons for details, duplicate,
delete, or more commands. These cells make dense screens efficient, but they can
also become the most fragile part of the UI if the implementation confuses a
renderer with a live component.

\begin{figure}[htbp]
\centering
\includegraphics[width=0.96\linewidth]{\practicalInteractiveCellsFigure}
\caption{Interactive cells need stable row identity and column identity. The
renderer paints badges and affordances, while editors, table actions, and popup
commands route through row ids, column keys, and Corusco command metadata.}
\label{fig:practical-interactive-cells}
\end{figure}

The first rule is the Swing table rule: a renderer is not an installed child
component. A renderer is used as a stamp. Swing asks it to configure a component
for a cell and then paints that component. The component returned by the
renderer does not receive mouse events as a normal child. If the renderer
contains a button, that button will look like a button but will not behave like a
normal button. This is why "button in table cell" examples that work by adding a
JButton to the renderer usually fail under keyboard, sorting, focus, and
accessibility review.

\begin{principlebox}{Renderers Show State; Editors Perform Interaction}
Use renderers to paint the current cell state. Use a table cell editor, table
input action, or explicit mouse-to-command router for interaction. Do not expect
renderer child components to behave as live controls.
\end{principlebox}

There are three implementation patterns. A real \api{TableCellEditor} is the
right choice when the user enters an editing mode, changes a value, and commits
or cancels that edit. A mouse-to-command router is the right choice when a click
on a painted affordance invokes a row command immediately. An overlay component
is the right choice when the interaction is complex enough to require a real
component but should appear visually over the cell. The overlay pattern is rare
and should be owned by the table view, not by a renderer pretending to be live.

\begin{tabularx}{0.96\linewidth}{@{}p{0.24\linewidth}X@{}}
\toprule
Pattern & Use when\\
\midrule
Renderer only & The cell is visual state: badge, progress, warning, formatted text.\\
Cell editor & The user edits a value with commit and cancel semantics.\\
Mouse router & A click invokes a command without entering an edit session.\\
Overlay component & A rich interaction must appear over the cell while the table remains visible.\\
\bottomrule
\end{tabularx}

Row identity is the second rule. A table may be sorted, filtered, grouped, or
partly refreshed while the user interacts. View row 5 is not a business row. It
is the row currently displayed at index 5. Every command must convert view
coordinates to model coordinates and then resolve a stable row id. If a command
stores only the view row, it may act on the wrong record after sorting or after a
background refresh. Corusco table descriptors are valuable here because the
column and row identity can be descriptor-backed rather than inferred from the
header text or renderer class.

\begin{lstlisting}[caption={Interactive cell commands convert coordinates before resolving row identity.}]
void invokeCellCommand(JTable table, Point point) {
    int viewRow = table.rowAtPoint(point);
    int viewColumn = table.columnAtPoint(point);
    if (viewRow < 0 || viewColumn < 0) {
        return;
    }

    int modelRow = table.convertRowIndexToModel(viewRow);
    int modelColumn = table.convertColumnIndexToModel(viewColumn);
    OrderRow row = model.rowAt(modelRow);
    ColumnKey column = descriptors.columnKey(modelColumn);

    commands.forCell(row.id(), column)
            .ifPresent(command -> command.invoke(CommandOrigin.mouse(point)));
}
\end{lstlisting}

This example is more important for what it refuses to do. It does not inspect the
localized header value. It does not use view row as a durable id. It does not
ask the renderer what command it thinks it represents. The command is resolved
from row identity and column identity. The renderer can paint the same command
state, but it is not the authority.

Hit testing is a design decision. Some interactive cells make the entire cell a
command target. A status cell may open details anywhere in the cell. Other cells
have multiple painted affordances, such as retry and cancel icons in the same
cell. In that case the renderer policy should expose a hit map calculation that
is independent of painting side effects. The same geometry method can be used by
the renderer, mouse router, tooltip provider, and tests.

\begin{lstlisting}[caption={A cell interaction policy computes hit zones from cell geometry.}]
record CellHitZone(String commandId, Rectangle bounds) {
    boolean contains(Point cellPoint) {
        return bounds.contains(cellPoint);
    }
}

interface InteractiveCellPolicy<R> {
    List<CellHitZone> hitZones(
            R row,
            ColumnKey column,
            Rectangle cellBounds,
            FontMetrics metrics,
            ComponentOrientation orientation);
}
\end{lstlisting}

The policy should work in cell coordinates, not screen coordinates. The table
mouse listener translates the event point into the cell rectangle, subtracts the
cell origin, and asks the policy which zone was hit. The renderer uses the same
cell bounds to paint icons and hover feedback. This removes a common class of
bugs where the painted icon and clickable icon drift apart when font size,
Look-and-Feel insets, row height, or component orientation changes.

Hover feedback needs table repaint discipline. A table does not repaint a cell
because an internal renderer button changed rollover state; there is no internal
button. Store hovered row, column, and hit zone in the table interaction model,
then repaint the old and new cell rectangles when the mouse moves. Keep hover
state in view coordinates only as a transient painting optimization, and clear
it when the model changes, sorter changes, table loses focus, or mouse exits.
Commands should still use row ids when invoked.

\begin{lstlisting}[caption={Hover state is transient view state; command state is model state.}]
final class TableInteractionState {
    private CellAddress hoveredCell;
    private String hoveredZone;

    void updateHover(JTable table, MouseEvent event) {
        CellAddress previous = hoveredCell;
        hoveredCell = CellAddress.fromViewPoint(table, event.getPoint());
        hoveredZone = hitZoneAt(table, hoveredCell, event.getPoint()).orElse(null);
        repaint(table, previous);
        repaint(table, hoveredCell);
    }
}
\end{lstlisting}

Keyboard operation should not be an afterthought. If a mouse click can invoke a
cell command, keyboard users need a route as well. The simplest route is to bind
Enter or Space on the selected cell to the default command for that column. If a
cell has multiple commands, the table can open a small context menu with the
cell commands, using the selected row and column as origin. Shift+F10 should open
the same menu. A row action column can also expose actions through the table's
main action map so accelerators work regardless of focus inside the table.

\begin{lstlisting}[caption={The table action map invokes the selected cell command.}]
table.getInputMap(JComponent.WHEN_ANCESTOR_OF_FOCUSED_COMPONENT)
        .put(KeyStroke.getKeyStroke("SPACE"), "invokeSelectedCell");
table.getActionMap().put("invokeSelectedCell", new AbstractAction() {
    @Override
    public void actionPerformed(ActionEvent event) {
        int viewRow = table.getSelectedRow();
        int viewColumn = table.getSelectedColumn();
        invokeDefaultCellCommand(table, viewRow, viewColumn, CommandOrigin.keyboard());
    }
});
\end{lstlisting}

Real cell editors are still the right tool for many cases. A tri-state checkbox
cell, a quantity stepper, a date picker, and an inline combo box need edit
sessions. The editor component receives focus, edits a value, and calls
\api{stopCellEditing()} or \api{cancelCellEditing()}. The editor must publish
the value through \api{getCellEditorValue()} and the table model must decide
whether the value is accepted. Validation errors should not be hidden in the
editor component. They should target the row and column so the renderer can show
the problem after editing ends.

\begin{lstlisting}[caption={A cell editor commits through the table model, not directly into a row object.}]
final class QuantityEditor extends AbstractCellEditor implements TableCellEditor {
    private final QuantityStepper stepper = new QuantityStepper();
    private Quantity draft;

    @Override
    public Component getTableCellEditorComponent(
            JTable table,
            Object value,
            boolean selected,
            int row,
            int column) {
        draft = (Quantity) value;
        stepper.setQuantity(draft);
        return stepper;
    }

    @Override
    public Object getCellEditorValue() {
        return stepper.quantity();
    }
}
\end{lstlisting}

This code is only half of the contract. The table model's \api{setValueAt()}
must validate the value and update row state through the presenter or row model.
If validation can fail, the editor may be allowed to remain open, or the edit may
commit a draft problem. Both policies can be correct. What is not correct is an
editor that mutates a domain object directly and leaves the table model unaware
that sorting, filtering, totals, or dirty state changed.

Sorting and filtering create subtle failures. Suppose the user clicks "Approve"
in a cell. The row becomes approved and the current filter hides approved rows.
Should the row disappear immediately? Maybe yes. But if it disappears before the
command result is announced, users may think nothing happened. A command handler
can decide to keep the row visible until refresh, show a notification, or move
selection to the next visible row. That is workflow logic. The cell renderer
should not make the decision by repainting itself as approved and hoping the
sorter catches up.

Popup menus inside cells should be routed like header popups. The popup origin is
the cell rectangle. The command list comes from row id and column key. Menu item
enabled state comes from command metadata. When the popup closes, transient
hover and pressed states are cleared. If the table is refreshed while the popup
is open, the command should either resolve the row id again or close the popup.
Never let a popup keep a stale row object that may have been replaced by the
model.

\begin{lstlisting}[caption={A row command popup is built from row id and column key.}]
void showCellMenu(JTable table, int viewRow, int viewColumn) {
    RowId rowId = rowIdAtViewRow(table, viewRow);
    ColumnKey column = columnKeyAtViewColumn(table, viewColumn);
    JPopupMenu menu = new JPopupMenu();

    for (Command command : commands.forCellMenu(rowId, column)) {
        JMenuItem item = new JMenuItem(command.toAction());
        item.setEnabled(command.enabled());
        item.setToolTipText(command.disabledReason().orElse(null));
        menu.add(item);
    }

    Rectangle rect = table.getCellRect(viewRow, viewColumn, true);
    menu.show(table, rect.x, rect.y + rect.height);
}
\end{lstlisting}

Corusco can make interactive cells more maintainable by moving column knowledge
out of renderer classes. A generated or descriptor-backed table column can carry
the column key, display resources, value type, alignment, formatter, problem
target, help topic, and command contributions. The interactive-cell installer
then asks the descriptor what the cell means. This is especially important for
multi-language applications. A column whose header is "Status" in English and
"Stav" in Czech should still have the same command behavior because behavior is
attached to the descriptor key, not to the rendered string.

\begin{lstlisting}[caption={Descriptor-backed columns provide command metadata to cells.}]
record InteractiveColumn<R>(
        ColumnKey key,
        ValueReader<R, ?> value,
        CellRendererPolicy<R> renderer,
        InteractiveCellPolicy<R> interaction,
        CommandContributor<R> commands) {
}
\end{lstlisting}

Accessibility has two levels. A rendered status badge should expose useful cell
text through the table model, renderer accessible context, or table accessible
implementation. An actual editor should expose its child component semantics
while editing. For immediate command cells, the table should expose the command
through keyboard and context menu, because assistive technology users cannot
click a painted icon that is not a real component. Tooltips should also be
available from keyboard focus, not only mouse hover.

Testing interactive cells is mostly testing coordinate and command contracts.
Unit tests can verify that the hit-zone policy returns the same zones under
left-to-right and right-to-left orientation, large fonts, and narrow columns.
Model tests can verify that a row command uses row id and column key. Swing tests
can select a cell and invoke Space, Shift+F10, and mouse click paths. Screenshot
tests can verify that hover, disabled, selected, focused, and error states remain
readable under the supported Look-and-Feels.

\begin{tabularx}{0.96\linewidth}{@{}p{0.24\linewidth}X@{}}
\toprule
Failure mode & Design response\\
\midrule
Button renderer does not click & Use a real editor or mouse-to-command router.\\
Wrong row after sorting & Convert view row to model row and resolve stable row id.\\
Clickable zone drifts & Share geometry between renderer, router, tooltip, and tests.\\
Keyboard cannot invoke action & Bind table actions and context menu to selected cell.\\
Popup acts on stale row & Resolve by row id or close popup on model refresh.\\
Localized header changes behavior & Attach behavior to column descriptors, not header text.\\
\bottomrule
\end{tabularx}

\begin{principlebox}{Lessons Learned From Interactive Cells}
Interactive table cells are table interactions first and painted widgets second.
The durable facts are row id, column key, command metadata, edit value, and
problem target. Corusco descriptors give those facts stable names, while Swing
renderers, editors, input maps, and popup menus provide the visible interaction
without pretending that renderer children are live components.
\end{principlebox}

\section{Additional Business Widgets}

\begin{roadmapbox}
This catalog section connects the larger component patterns to smaller controls
that appear on business screens every day. The goal is not to make each one a
full framework. The goal is to identify the value, command, focus, validation,
and accessibility contract that keeps even compact widgets from becoming hidden
state in a layout.
\end{roadmapbox}

The required controls cover the deep cases, but real business screens contain several
other recurring non-stock widgets. They do not all deserve a full chapter. They do
deserve a construction rule because the same mistakes repeat: state hidden in child
components, color-only communication, focus paths that work only by mouse, and commands
that cannot be tested without a screen.

\begin{figure}[htbp]
\centering
\includegraphics[width=0.96\linewidth]{\practicalBusinessWidgetsFigure}
\caption{Small business widgets also need named state. Token fields, status badges, tri-state controls, step panels, and overflow command bars become maintainable when their values and commands are not hidden in child components.}
\label{fig:practical-business-widgets}
\end{figure}

\begin{principlebox}{Lessons Learned From the Business-Widget Catalog}
Small widgets fail for the same reasons large widgets fail: unclear value
ownership, mouse-only commands, missing problem targets, and visual state that is
not represented in the model. A compact control still deserves a compact
contract.
\end{principlebox}

\section{Search Fields With Scope Selectors}

\begin{roadmapbox}
This section treats scoped search as one query model with command state. It
names the scope, query, parsing policy, clear operation, submit operation, and
validation response so a search field can be reused without embedding repository
calls in its Swing listeners.
\end{roadmapbox}

A search field with a scope selector combines query text, search command, clear command,
and scope choice. It may appear as \texttt{[ customer id v ][ query ][ Search ]} or as
one visual field with a left scope segment and right clear icon. The scope changes
parsing and search semantics, so it belongs in the model. It is not decoration.

\begin{lstlisting}[caption={Scoped search state keeps query and scope together.}]
record SearchDraft(
        SearchScope scope,
        String query,
        boolean caseSensitive,
        List<Problem> problems) {

    boolean canSubmit() {
        return !query.isBlank() && problems.isEmpty();
    }
}
\end{lstlisting}

The search action reads a draft, validates it, and submits a request. The scope menu
changes draft state. The clear command clears query and perhaps resets problems. If
changing scope invalidates the query, the problem model should say so. Do not let the
search field call a repository directly just because the Search button is nearby. The
field is an editor and command surface; the presenter owns the workflow.

\begin{principlebox}{Lessons Learned From Scoped Search}
The scope is part of the query meaning, not a visual prefix. The search action
should read a draft that contains both scope and text, then let the presenter
decide whether the query can run. That keeps parsing, validation, permissions,
and async search behavior out of the text-field component.
\end{principlebox}

\section{Token and Chip Fields}

\begin{roadmapbox}
This section introduces token fields as list editors with a text-entry tail. It
focuses on selected-token identity, query text, suggestions, focused token,
removal commands, and keyboard behavior. The deeper tag-cloud section expands
the same ideas for catalog-backed tags.
\end{roadmapbox}

Token fields represent a list of selected items such as tags, recipients, products,
roles, or filters. The user can type to search, choose suggestions, remove tokens, and
sometimes reorder them. A token field is not just a text field with painted labels. It
is a list editor with a text input. The model should contain tokens, current query,
suggestion state, focused token, and problems.

Keyboard behavior is the main design challenge. Backspace in an empty query may focus
the previous token, then remove it on a second press. Arrow keys may move between tokens
and the query. Enter may choose the highlighted suggestion. Escape may close
suggestions. Mouse behavior must match the same model. A token field that works only
with a mouse will frustrate fast data entry users.

\begin{tabularx}{0.96\linewidth}{@{}p{0.25\linewidth}X@{}}
\toprule
Token field fact & Owner\\
\midrule
Committed tokens & Token list model or presenter state.\\
Current query & Text document or query model.\\
Suggestions & Suggestion model, often loaded asynchronously.\\
Focused token & Field focus model, not a painted rectangle only.\\
Validation & Problem model targeting token list or specific token id.\\
\bottomrule
\end{tabularx}

\begin{principlebox}{Lessons Learned From Token Fields}
A token field is a list editor first and a text field second. Treating the chips
as text decoration loses focus, identity, duplicate policy, and validation. A
model that names those facts makes removal, suggestions, and keyboard navigation
reviewable.
\end{principlebox}

\section{Tri-State and Nullable Boolean Controls}

\begin{roadmapbox}
This section handles the common case where Yes and No are not enough. It chooses
an explicit nullable or tri-state model, describes cycling and menu policies,
and explains how to communicate an indeterminate state without overloading
enabled, selected, or disabled flags.
\end{roadmapbox}

Business forms often distinguish Yes, No, and Unknown. A normal checkbox is binary. A
nullable Boolean, inherited setting, or mixed selection needs a tri-state control. The
control may look like a checkbox with an indeterminate state, a segmented control, or a
combo-like selector. The model should name the three values rather than overloading
selected and enabled flags.

\begin{lstlisting}[caption={Tri-state values should be explicit.}]
enum TriState {
    YES,
    NO,
    UNKNOWN
}
\end{lstlisting}

The cycling order should be deliberate. Some controls cycle Unknown, Yes, No. Others
cycle Yes, No and expose Unknown only through a menu. Accessibility must announce the
state. Validation must know whether Unknown is allowed. Persisted values should not
confuse null with false. A tri-state checkbox is small, but it is a business semantics
control, not a visual novelty.

\begin{principlebox}{Lessons Learned From Tri-State Controls}
Unknown is a business value, not a broken checkbox. Once the third state is
explicit, validation, persistence, accessibility text, and bulk-edit mixed state
can all agree on what the user is seeing and committing.
\end{principlebox}

\section{Status Badges and Severity Pills}

\begin{roadmapbox}
This section designs compact state indicators that remain accessible. It covers
severity identity, label resources, color and icon policy, tooltip text,
renderer reuse, and the rule that status must never be communicated only through
color.
\end{roadmapbox}

Status badges are compact labels such as Active, Suspended, Draft, Overdue, Failed,
Warning, or Synced. They often use color. Color is useful but not sufficient. The badge
text, accessible name, tooltip, and exported value should carry the same status. A
renderer or component should receive a status object with severity, label key, optional
explanation, and maybe help topic. It should not infer severity from a localized string.

\begin{lstlisting}[caption={A badge receives status metadata, not only text.}]
record StatusBadgeModel(
        Severity severity,
        String labelKey,
        Optional<String> explanationKey) {
}
\end{lstlisting}

Badges should be cheap to render because they often appear in tables. Avoid allocating
new fonts, borders, and colors per cell. Cache derived presentation objects by severity
and Look-and-Feel generation. In a form, a badge can be a real component. In a table, it
is usually renderer state. The same model can feed both.

\begin{principlebox}{Lessons Learned From Status Badges}
A badge is a renderer for a named state. Color may support that state, but text,
icon, tooltip, and accessible description must also carry it. Corusco resources
and problem metadata keep those representations synchronized across forms and
tables.
\end{principlebox}

\section{Inline Validation Summary Strips}

\begin{roadmapbox}
This section turns a list of validation problems into a compact screen control.
It decides what problems are summarized, how users navigate from summary to
field, when the strip appears, and how it avoids becoming a second inconsistent
validation system.
\end{roadmapbox}

An inline validation summary strip appears above or below a form region and summarizes
current problems. It may show "3 problems", the first problem, or grouped messages. The
strip should observe the problem model. It should not scrape red borders or tooltips. It
should route clicks or keyboard commands to problem targets when possible.

The summary is also a recovery tool. If the first error is hidden inside a collapsed
section or tab, the summary can navigate there. If a table cell has a row and column
target, the summary can select and scroll to it. If a problem comes from asynchronous
validation, the summary can show pending state. This requires structured problem
targets. Text-only problem lists are less useful.

\begin{principlebox}{Lessons Learned From Validation Strips}
A validation strip should be a view of the same problem model used by the fields.
Its value is navigation and prioritization, not duplicate validation logic. When
it can focus the exact problem target, it becomes useful instead of ornamental.
\end{principlebox}

\section{Step and Wizard Panels}

\begin{roadmapbox}
This section treats a step panel as a state machine over staged workflow. It
names steps, completion, skipped state, blocking problems, command availability,
and navigation rules so the panel can represent real workflow rather than just a
row of labels.
\end{roadmapbox}

Step panels are common for imports, setup tasks, batch operations, and approval
workflows. They are often implemented as a card layout plus Back and Next buttons. The
practical design needs a step model: current step, step validity, visited state, command
availability, summary, and optional task state. The visible step indicator should
observe the model. It should not be the model.

\begin{lstlisting}[caption={A wizard step model names navigation state.}]
record StepState(
        String stepId,
        boolean current,
        boolean complete,
        boolean hasProblems,
        boolean visited) {
}
\end{lstlisting}

Navigation commands should validate the current step before moving when the workflow
requires it. Back may be allowed even when the current step is invalid. Cancel may
require confirmation if any step is dirty. Finish may start a background task. A wizard
is a screen workflow with a compact visual header, not a collection of cards with
buttons hard-coded on each card.

\begin{principlebox}{Lessons Learned From Step Panels}
Steps are workflow state, not layout decoration. A good step panel can explain
which step is current, which steps are blocked, what command moves next, and
which validation problems prevent movement. That requires a model richer than a
selected index.
\end{principlebox}

\section{Compact Command Bars With Overflow}

\begin{roadmapbox}
This section designs command bars that keep primary commands visible and move
secondary commands into overflow. It covers command priority, responsive
measurement, keyboard access, popup rebuilding, and the need for the same action
object to serve toolbar and menu presentations.
\end{roadmapbox}

Dense business screens often have more commands than fit in a toolbar. A compact command
bar shows primary commands and moves secondary commands into an overflow menu. The
command bar should be generated from command metadata or a presenter command list. It
should not be laid out by hand in several places. Overflow is a presentation decision;
command identity and enablement remain stable.

The overflow menu should preserve keyboard shortcuts, disabled reasons, and selection
state. If a command moves from visible button to overflow due to width, tests should
still locate it by command id. Accessibility should not lose the command because it
moved. Screenshot tests should include narrow widths so overflow behavior is exercised.

\begin{principlebox}{Lessons Learned From Command Bars}
Overflow is a presentation choice, not a second command model. When toolbar
buttons and overflow menu items share generated command metadata, resizing the
bar changes only placement. Permissions, labels, shortcuts, help, and disabled
reasons remain one source of truth.
\end{principlebox}

\section{Design walk: small widgets still need full contracts}

\begin{roadmapbox}
This walk compares the small catalog widgets through one lens: what goes wrong
when a compact component skips model design. It extracts a reusable checklist
for value identity, command ownership, event routing, problem targeting,
Look-and-Feel respect, and cleanup.
\end{roadmapbox}

The catalog widgets are smaller than calendars or table headers, but they are not exempt
from design. A token field is a list editor. A status badge is a semantic renderer. A
tri-state control is a value editor with three business states. A validation strip is a
navigator over problem targets. A step panel is a workflow map. A command overflow bar
is a responsive command presenter. Each can be implemented with a few Swing components,
but each can also become unmaintainable if its state lives in labels and listeners.

Start with the question "what would Corusco key?" A token field should have a field key
for the whole token list and perhaps component keys for the query editor and suggestion
popup. A status badge should have a resource key for visible text and a stable severity
or status code. A tri-state selector should have a field key and explicit enum value. A
validation strip should consume \api{ProblemSet} and
\api{ProblemTarget}. A step panel
should have step ids and command keys. An overflow bar should use action keys. If the
answer is "nothing can be keyed because the widget is just visual", the widget is
probably hiding application meaning.

\begin{lstlisting}[caption={Catalog widgets should expose semantic state, not only children.}]
interface TokenFieldView<T> {
    List<T> tokens();
    String query();
    Optional<T> focusedToken();
    void setProblems(ProblemSet problems);
}

interface StepPanelView {
    void setSteps(List<StepState> steps);
    void setCurrentStep(String stepId);
}
\end{lstlisting}

These interfaces are intentionally small. They say what a binding needs to know without
prescribing painting. A token field implementation may use labels and a text field,
custom painting, or a third-party chip component. A step panel may use buttons, labels,
or a custom timeline. The Corusco binding cares about tokens, query, problems, steps,
and current step. This lets a screen migrate implementation style without changing
validation, commands, or tests.

Token fields have the richest event path among the catalog widgets. Typing changes the
query. Query changes may start asynchronous suggestions. Suggestion completion must be
generation-checked so old results do not replace new ones. Clicking or pressing Enter on
a suggestion adds a token and clears or keeps the query according to product policy.
Backspace in an empty query focuses or removes a token. Removing a token updates
validation. A duplicate token may be rejected, merged, or allowed depending on the
field. None of this belongs inside a chip label. The chip label should render a token
and expose remove action. The model owns the token list and query workflow.

\begin{lstlisting}[caption={Token field events update list state through named operations.}]
final class TokenFieldModel<T> {
    void setQuery(String query) { /* update draft query */ }
    void acceptSuggestion(T token) { /* add token by policy */ }
    void removeToken(T token) { /* remove token and validate */ }
    void focusPreviousToken() { /* keyboard navigation */ }
    void clearFocusedToken() { /* delete selected token */ }
}
\end{lstlisting}

Corusco observable lists fit naturally here. The token list can be an
\api{ObservableList}. A suggestion list can be a loadable row source. The
field binding can adapt list changes to chip components. A problem model can target the
token field or a specific token id. A task service can load suggestions away from the
EDT and return them through a generation counter. The visual control remains a compact
field, but the behavior is assembled from named Corusco concepts instead of hidden
state.

Status badges have a different risk: they look so simple that developers often pass only
a string and color. That breaks localization, accessibility, export, filtering, and
testing. A status badge should receive status metadata: stable status code, severity,
label resource key, optional explanation, optional help topic, and optional icon key. In
a table renderer, the badge is presentation of a row value. In a form, it may be a
component bound to state. Either way, color is an output, not the model.

\begin{lstlisting}[caption={Status badge policy derives visuals from semantic state.}]
record StatusPresentation(
        String statusCode,
        ProblemSeverity severity,
        ResourceKey labelKey,
        Optional<ResourceKey> explanationKey,
        Optional<HelpTopic> helpTopic) {
}
\end{lstlisting}

The component wiring is then predictable. The presenter publishes a status code or
status value. A policy maps it to \api{StatusPresentation}. Resources resolve text. The
badge renderer chooses colors from severity and UI defaults. Accessibility receives
label and explanation. Export uses the status value or resource text according to export
policy. Tests assert status code and presentation mapping separately from screenshots. A
screenshot proves the badge is legible; it is not the only proof of status semantics.

Tri-state controls need a value model because Swing's stock checkbox is binary. Some
Look-and-Feels expose indeterminate state through client properties, but the application
should still model the value as Yes, No, and Unknown or Inherited. The cycling order is
a product decision. For a filter, Unknown may mean "any"; for a form, Unknown may mean
"not supplied"; for a bulk editor, Mixed may mean "selected records differ." These are
not the same state. Use an enum that names the business meaning.

The event path for a tri-state control should distinguish user cycling from programmatic
load. Loading a mixed selection sets Mixed without marking the field dirty. User click
cycles to Yes or No and marks dirty. Reset returns to the loaded state. Validation may
reject Unknown for required fields. Corusco field models can represent this as a
combo-like value with a custom renderer, or as a custom component adapter. What matters
is that null is not silently treated as false.

Validation summary strips are often underestimated. A useful strip is not a static label
with red text. It observes \api{ProblemSet}, groups or orders problems, exposes a
primary recovery action, and navigates to targets. If the first problem is inside a tab,
the strip can activate the tab. If the problem targets a table cell, the strip can
select and scroll to the row and column. If the problem targets a buddy field, the strip
can focus the text child or wrapper according to field policy. That is only possible
when problems have structured targets and the view has component keys.

\begin{lstlisting}[caption={A validation strip routes by problem target.}]
void activateProblem(Problem problem) {
    ProblemTarget target = problem.target();
    ComponentKey componentKey = focusResolver.resolve(target);
    JComponent component = componentRegistry.require(componentKey);
    focusRepair.focus(component);
}
\end{lstlisting}

The strip does not know the component tree by accident. It uses a resolver that maps
Corusco problem targets to component keys. This keeps the summary usable with custom
fields, table cells, nested panels, and dialogs. It also keeps tests stable: a test can
assert that a problem for
\api{OrderFields.SHIPPING\_WEIGHT} focuses the measurement field, not whatever
component happens to have the current English label "Weight".

Step panels and wizard panels are workflow widgets. The model should state which step is
current, which steps are complete, which are blocked, which have problems, and which
commands are available. The view should not decide that a step is complete because its
label is green. The presenter or workflow model decides completion. Commands such as
Back, Next, Finish, Cancel, and Help are normal commands, preferably with generated or
typed action keys. The step panel is then a visualization of workflow state plus a
command surface.

\begin{lstlisting}[caption={Wizard commands read workflow state, not card layout.}]
record WizardState(
        String currentStepId,
        List<StepState> steps,
        boolean canGoBack,
        boolean canGoNext,
        boolean canFinish,
        ProblemSet problems) {
}
\end{lstlisting}

The event path for Next is important. It may commit active editors in the current step,
run local validation, maybe run asynchronous validation, then move to the next step only
if the current step can be left. Back may skip validation. Finish may start a task and
make the whole wizard busy. Cancel may ask for dirty confirmation. None of these
policies should live in the visual step indicator. The step indicator observes
\api{WizardState}. Command handlers own workflow.

Overflow command bars are also model-driven. The command list is stable: Save, Refresh,
Export, Print, Delete, Help. The available width changes where commands are presented.
Some commands are visible as buttons; others move to the overflow menu. The command id,
enabled state, accelerator, help topic, and disabled reason do not change. Corusco
command metadata is a good fit because it decouples command identity from the current
surface. The same
\api{ActionKey} can be rendered as toolbar button or menu item.

Testing catalog widgets follows their model. Token field tests should cover query,
suggestion generation, duplicate handling, removal, keyboard deletion, and validation.
Badge tests should cover status mapping and accessible text. Tri-state tests should
cover load, cycle, reset, and required validation. Validation strip tests should cover
problem ordering and focus routing. Step panel tests should cover command availability
and navigation policy. Overflow bar tests should cover command identity at wide and
narrow widths. Screenshots then prove the assembled controls remain legible and
pleasant.

\begin{tabularx}{0.96\linewidth}{@{}p{0.22\linewidth}p{0.35\linewidth}X@{}}
\toprule
Widget & Model decision & Corusco help\\
\midrule
Token field & Are duplicates allowed, and how does suggestion loading cancel? & Observable lists, task service, generation counter, problem targets.\\
Status badge & Is the state informational, warning, error, or workflow status? & Resource keys, problem severity, help topics, renderer policies.\\
Tri-state & Does the third value mean unknown, inherited, mixed, or any? & Typed field value, validation rules, generated resources.\\
Validation strip & Which problem is primary and where does recovery focus? & ProblemSet, ProblemTarget, component keys, focus resolver.\\
Step panel & Which steps can be left, skipped, revisited, or completed? & Command keys, problem models, lifecycle scopes, task state.\\
Overflow bar & Which commands remain primary at each width? & Command descriptors, SwingActionAdapter, resource-backed menu items.\\
\bottomrule
\end{tabularx}

\begin{principlebox}{Lessons Learned From the Small-Widget Walk}
The smaller the widget, the more tempting it is to hide state in the visual
surface. That is exactly where Corusco-style keys and command descriptors help:
they force the component to expose the value, action, and problem facts that a
reviewer needs.
\end{principlebox}

\section{Worked Construction: An Order Filter Panel}

\begin{roadmapbox}
This worked example assembles several earlier ideas into one realistic filter
panel. It starts with a domain model, maps each field to a control contract,
wires commands and validation, explains what Corusco owns, and shows how the
panel can remain a normal Swing surface rather than a custom mini-framework.
\end{roadmapbox}

A practical way to read this chapter is to imagine one dense business panel.
The panel filters orders by customer, status, weight, due date, and date-time
span. It has a split Export button, a customer search field with Clear and
Search buddies, a weight field with a unit selector, a table with descriptor
backed headers, a date-time span selector, a validation summary strip, and a
compact command bar. This is not a large screen, but it is large enough to
force every design decision that small examples usually avoid.

Begin with the facts, not the components. The committed filter value might be
an \api{OrderFilter}. The draft state contains raw query text, selected scope,
selected statuses, measurement draft, temporal draft, sort state, and problem
state. The commands are Search, Clear, Export PDF, Export CSV, Export XLSX,
Reset Dates, Today, This Week, and Help. The table columns have stable ids.
The validation problems have field keys and problem codes. The view has
component keys for the fields, table, header, summary, and command surfaces.
Only after those facts exist should the Swing panel be assembled.

\begin{lstlisting}[caption={A practical panel starts with a draft model.}]
record OrderFilterDraft(
        SearchDraft customerSearch,
        Set<OrderStatus> statuses,
        MeasurementDraft minimumWeight,
        DateTimeSpanDraft dueWindow,
        SortState sortState,
        ProblemSet problems) {

    boolean canSearch() {
        return problems.isEmpty() && !customerSearch.query().isBlank();
    }
}
\end{lstlisting}

The draft is not the committed search request. It is the editable state of the
panel. It can contain invalid amount text. It can contain a date range whose
end is before its start. It can contain a customer query whose scope requires
digits but currently has letters. A committed request should be produced only
when conversion and validation succeed. Corusco form models use the same
separation, so custom widgets should cooperate with it instead of bypassing
it.

The view should expose logical component slots. It may use MigLayout, nested
panels, custom field wrappers, and a table. The presenter and bindings should
not need to know the exact layout. They should need only the text editor for
customer query, the search scope selector, the measurement field view, the
date-time span view, the table, the validation summary, and command surfaces.
This keeps layout change from becoming a binding rewrite.

\begin{lstlisting}[caption={The view contract names logical slots, not layout details.}]
interface OrderFilterView {
    SearchFieldView customerSearch();
    MeasurementFieldView minimumWeight();
    DateTimeSpanView dueWindow();
    JTable ordersTable();
    ValidationSummaryView validationSummary();
    CommandSurface commands();
}
\end{lstlisting}

A generated view contract could contain similar names when the fields are
known to the generator. A handwritten view can still follow the same rule.
The interface says what the screen needs to bind. It does not say that the
customer field sits in row one, column two. It does not say that the unit
button is the third child of a panel. That distinction is where maintainable
Swing starts: layout is visible, but meaning is not trapped in layout.

Next decide which state is local to the panel and which belongs to a presenter.
Draft editing state can live in a form model owned by the presenter. The table
row source may live in a loadable list. Table state may be owned by
\api{TableStateController}. The command set belongs to the presenter because
commands express workflow. The view owns component construction and direct
Swing lifecycle. The binding scope owns installed relationships. This
ownership map should be written before code, even if only in a review note.

\begin{tabularx}{0.96\linewidth}{@{}p{0.27\linewidth}X@{}}
\toprule
Fact & Owner in the order filter panel\\
\midrule
Raw customer query & Search field model or generated text field model.\\
Search scope & Search draft model, not the scope button label.\\
Weight amount and unit & Measurement field model with converter and unit policy.\\
Due date span & Temporal field model with zone and endpoint policy.\\
Visible order rows & Observable row source or loadable list.\\
Column widths and visibility & Table state controller by table and column ids.\\
Export commands & Presenter command set adapted to Swing actions.\\
Validation summary & Problem model plus focus resolver.\\
Listeners and bindings & View activation scope.\\
\bottomrule
\end{tabularx}

The construction order follows ownership. Build the model. Build the command
set. Build the view. Install editor bindings. Install command bindings.
Install table descriptor, sorter, state controller, and header behavior.
Install validation decoration and summary routing. Install accessibility
metadata. Register all returned bindings and controllers in the activation
scope. Finally, render screenshots of the normal, invalid, busy, and narrow
states. This order is slower to write than a quick prototype, but it is much
faster to review six months later.

\begin{lstlisting}[caption={A practical assembly reads as a sequence of owners.}]
BindingScope scope = new BindingScope();
OrderFilterModel model = new OrderFilterModel(clock, resources);
OrderFilterView view = new OrderFilterPanel(resources);
OrderFilterCommands commands = new OrderFilterCommands(model, services);

scope.add(bindSearchField(view.customerSearch(), model.customerSearch()));
scope.add(bindMeasurement(view.minimumWeight(), model.minimumWeight()));
scope.add(bindDateTimeSpan(view.dueWindow(), model.dueWindow()));
scope.add(bindCommands(view.commands(), commands));
scope.add(bindProblems(view.validationSummary(), model.problems()));
scope.add(installOrderTable(view.ordersTable(), model.rows(), resources));
\end{lstlisting}

The code is schematic, but the review value is concrete. Each line installs a
relationship with a name. Each relationship can be tested and closed. If a
future feature adds an "Include archived orders" tri-state control, it should
add a model fact, a view slot, a binding, a command or validation rule if
needed, and a screenshot state. It should not attach an anonymous listener to
the nearest checkbox and update the query string directly.

The event path for Search illustrates the whole design. The user types in the
customer search field. The document binding writes raw query to the search
model. The search model updates problems and command state. The Search action
observes command state and becomes enabled. The user presses Enter or clicks
Search. The command asks the filter model for a committed request. If commit
fails, problems are published and focus moves to the first problem. If commit
succeeds, the presenter submits a task. Busy state disables Search and export
commands. Task completion replaces rows on the EDT if the generation is still
current. The table model fires precise events. The selection and column state
remain stable.

The event path for Export is different. Export should not read visible table
cells or localized headers. It should read descriptor-backed rows and columns,
apply current filter and selection policy, resolve export headers through
resources, and run outside the EDT if output is slow. The split button is only
the command surface. Corusco table descriptors and command metadata do the
work of keeping export identity stable. A screenshot of the split button is
useful, but the export test should assert descriptor ids, not pixels.

The event path for date changes is different again. The user opens the date
span panel, chooses Today, edits the end time, and presses Apply. The temporal
model updates draft state for each action. Apply validates ordering, bounds,
zone conversion, and endpoint policy. If valid, it commits the span and marks
the filter dirty. If invalid, it publishes problems without closing. Cancel
restores committed state. Search reads only committed filter state unless the
product says Apply is implicit. These decisions should be visible in method
names and tests.

This worked panel is also where professional screenshots matter. A screenshot
should show the actual visual result of these decisions: the split button
looks like one command group; the unit field looks like one field; the table
header menu looks like a real header command surface; the temporal selector
shows Apply and Cancel; validation is visible without destroying layout. If
the screenshot is merely a schematic, it cannot prove the component design is
usable. If it is generated from deterministic Swing code, it can be reviewed
like any other artifact.

\begin{principlebox}{Lessons Learned From the Order Filter Panel}
The value of a compound panel is not that it hides all child controls. It is that
it exposes a coherent draft, command set, validation surface, and lifecycle
scope. Corusco makes the panel cheaper to maintain by supplying field and command
identity, while Swing remains responsible for the visible editing experience.
\end{principlebox}

\section{Model Design for Compound Components}

\begin{roadmapbox}
This section abstracts the model decisions repeated across the chapter. It
compares raw values, parsed values, draft values, committed values, selection,
focus, async state, and problems. The goal is to give readers a vocabulary for
choosing the smallest model that can still explain the component's behavior.
\end{roadmapbox}

Compound component models should be small, but not vague. They should avoid
both extremes: a model that is just a map of strings and a model that mirrors
every child component property. The model should contain facts that matter to
application behavior. For a unit field, those facts are raw amount, unit,
parse state, problems, and touched or dirty state. For a split button, they
are command choices and popup availability. For a calendar, they are visible
month, focused date, draft selection, committed value, and policy. If a fact
does not affect behavior, validation, accessibility, or persistence, it may
belong in the view instead.

The model should name incomplete input. Swing text components are good at
holding incomplete input. Domain records are usually not. A decimal field
cannot turn every intermediate text into \api{BigDecimal}. A date-time field
cannot turn every draft into \api{Instant}. A token field cannot turn every
query into a selected token. Corusco field models already acknowledge this by
separating raw text, parse state, and committed value. Custom component
models should do the same rather than forcing invalid input out of existence.

\begin{lstlisting}[caption={A custom model should preserve incomplete input.}]
record FieldDraft<T>(
        String rawText,
        Optional<T> parsedValue,
        ProblemSet problems,
        boolean touched,
        boolean dirty) {

    boolean committable() {
        return parsedValue.isPresent() && problems.isEmpty();
    }
}
\end{lstlisting}

This generic draft is not a proposed Corusco API. It is a design reminder.
Raw input, parsed value, problems, touched state, and dirty state are
different facts. Many custom components fail because they collapse them. A
unit field stores only \api{BigDecimal}; the user cannot type "-". A date
field stores only \api{LocalDate}; the user cannot type a partial date. A
token field stores only selected tokens; the current query is lost. A wizard
stores only current card index; validation state is inferred from button
enabled state. These shortcuts make demos shorter and applications weaker.

The model should also name policy. Unit conversion policy, date endpoint
policy, duplicate token policy, quick-range evaluation policy, disabled-date
policy, and split-button enablement policy are not implementation details.
They are product decisions that influence behavior and tests. A model that
contains values but not policy forces policy into listeners. A listener is a
poor place for policy because it is hard to reuse, hard to test, and easy to
skip from another input path.

\begin{tabularx}{0.96\linewidth}{@{}p{0.28\linewidth}X@{}}
\toprule
Policy & Why it belongs near the model\\
\midrule
Unit conversion & Changing the unit may preserve text, convert quantity, clear text, or reject.\\
Date endpoints & Storage, labels, and validation must agree on inclusive or exclusive end.\\
Duplicate tokens & Accept, reject, merge, or focus existing token is a business decision.\\
Split enablement & Main action and popup surface may have different availability.\\
Quick range capture & Relative ranges may evaluate when selected or when submitted.\\
Header visibility & Hideability belongs to descriptor capability, not menu layout.\\
\bottomrule
\end{tabularx}

Observable models should avoid over-notification. A unit field does not need
to publish a full screen refresh for every caret move unless caret state is a
model fact. A calendar does need to repaint focus when the focused day moves.
A token field does need to update suggestions when query changes, but it may
debounce service-backed suggestions. Corusco observable values and lists make
state changes explicit, but they do not remove the need to choose event
granularity. Too many events can make a screen feel unstable; too few events
can leave commands stale.

The model should support reset and close. Reset restores committed state or
loaded baseline. Close stops subscriptions, cancels suggestion loads, closes
table-state controllers, or invalidates generations. A view may be removed
from the screen while a model lives longer, or a model may close with the view.
The lifecycle must be explicit. Corusco scopes are useful because they give
the construction code a standard place to put these relationships. If a
custom widget starts a timer or listens to a model without returning a handle,
it is already making lifecycle harder.

\begin{principlebox}{Lessons Learned From Compound Models}
The best model is not the largest model. It is the model that names every fact
whose lifetime exceeds a single paint call or event callback. Draft text,
selection, command enablement, async loading, and validation all deserve names
when another part of the application must reason about them.
\end{principlebox}

\section{Event Wiring as a Reviewable Path}

\begin{roadmapbox}
This section follows events from Swing gesture to domain effect. It shows how
input maps, actions, document listeners, popup listeners, and model
subscriptions form a reviewable path. It also explains where Corusco command and
field bindings should sit so the view does not become a workflow controller.
\end{roadmapbox}

Every compound widget should have an event path that can be described in a
few sentences. The description should include the Swing event, the component
operation, the model mutation, the command or validation consequence, and the
view refresh. If the description jumps directly from "user clicked" to "the
screen updates somehow", there is hidden design work. Good Swing code makes
the path ordinary enough to review.

For text input, the path normally starts in the document. The document event
arrives on the EDT after the text mutation. The binding reads the document
text and writes raw text to the field model. The field model parses or
schedules parsing. Validation updates problems. Command state updates. The
view decoration observes problems and repaints. A formatter may run on focus
lost or commit, not on every keystroke unless explicitly designed. This path
should be the same whether the text field is plain, has a buddy button, or is
inside a unit field.

For popup input, the path starts with a mouse or key action. The action
resolves context, refreshes menu state, and shows the popup. The popup item
invokes a command or model operation. The model updates. The popup closes.
The component refreshes. If the popup item stores a view index, visible label,
or component reference that can become stale, the path is fragile. Resolve
context to stable keys before creating actions, or re-resolve deliberately at
invocation time.

For table header input, the path starts with pointer or focused-column
context. The header installer resolves descriptor identity. The action mutates
sorter, filter, visibility, help, or state. The renderer observes state. The
table-state controller persists if needed. The body selection model should
not be touched unless the command explicitly says it operates on rows. This
path is a good test of whether a team understands table coordinate systems.

For calendar input, the path starts with focus or pointer navigation. The
calendar model changes visible month, focused date, draft date, or committed
date according to the action. Validation checks disabled dates and bounds.
Apply commits. Cancel restores. The component repaints only the affected
cells when possible. A custom-painted calendar that handles mouse clicks but
not keyboard actions has only half an event path.

\begin{lstlisting}[caption={Event paths should converge on model operations.}]
void installMeasurementEvents(MeasurementFieldView view,
        MeasurementFieldModel model) {
    view.addAmountChangeListener(event ->
            model.setRawAmount(view.rawAmount()));
    view.addUnitSelectionListener(event ->
            model.selectUnit(view.unit()));
    view.addCommitListener(event ->
            model.commitDraft());
    view.addCancelListener(event ->
            model.resetDraft());
}
\end{lstlisting}

The listeners are small because the model operations have names. If a listener
contains parsing, conversion, validation, command enablement, tooltip text,
and border changes, it is not an event adapter. It is a hidden model. Hidden
models are difficult to share with keyboard paths, tests, accessibility
actions, and programmatic updates. Corusco helps because it gives bindings and
commands a place to live outside listeners.

Programmatic updates need their own path. Loading a record into a form should
not mark every field dirty. Resetting a field should not fire a search. A
model refresh should not open a popup. A table state restore should not be
persisted as a new user change before restore completes. Many Swing bugs come
from treating programmatic updates as if they were user input. Good models
can distinguish load, user edit, commit, reset, and close.

Reentrancy is a less common but important Swing aspect. Showing a popup,
dialog, file chooser, drag-and-drop interaction, or native menu may allow
other EDT work to run before the original method resumes. A listener that
captures state, shows UI, and then continues as if nothing changed may commit
stale data. For compound components, this means popup and chooser actions
should re-check lifecycle and current model state before committing results.
Corusco generation counters and scopes give a straightforward vocabulary for
rejecting stale outcomes.

\begin{principlebox}{Lessons Learned From Event Wiring}
An event path is maintainable when every hop changes vocabulary intentionally:
gesture, action, command, model update, binding refresh, repaint. Bugs cluster
where those translations are skipped and a mouse listener reaches directly into
domain state.
\end{principlebox}

\section{Look-and-Feel and Professional Visual Design}

\begin{roadmapbox}
This section turns the screenshots requirement into engineering rules. It covers
UI defaults, spacing, borders, density, disabled and error states, font metrics,
dark and light themes, and the difference between professional integration and a
custom component that only looks correct on one machine.
\end{roadmapbox}

A professional-looking Swing component usually starts by doing less custom
painting, not more. Use the current Look-and-Feel for text fields, buttons,
tables, menus, focus, disabled state, and font metrics. Compose existing
components with careful spacing and shared borders. Add custom painting only
where the visual concept cannot be expressed with stock components. The goal
is not to make every widget visually unique. The goal is to make custom
behavior feel native to the surrounding application.

The first aesthetic decision is density. Business applications need compact
screens, but compact does not mean cramped. Fields need enough horizontal
space for realistic values. Buttons need readable labels. Validation text
needs wrapping room. Tables need row height that supports status markers and
focus. A screenshot made with short English labels and tiny data is not a
professional proof. It should include realistic customer names, long paths,
currency, dates, status, disabled state, and at least one validation message.

The second aesthetic decision is grouping. A split button should look like one
command group. A unit field should look like one field. A validation strip
should belong to the form it summarizes. A table header menu should visually
belong to the table header. A date-time span selector should make start and
end feel paired. MigLayout is useful because it lets the code express these
relationships directly: aligned labels, growing fields, fixed command
segments, baseline alignment, and consistent gaps.

The third aesthetic decision is state. Normal, hover, focused, pressed,
disabled, invalid, busy, selected, and read-only states all need to be legible.
Do not remove focus painting to make a control look cleaner. Do not make a
disabled unit selector indistinguishable from a suffix label unless it is
truly read-only. Do not show validation only as a red border if the field also
has an explanation. Professional screenshots should capture meaningful states,
not only the happy path.

The fourth aesthetic decision is color restraint. Use color to communicate
state, not to decorate every surface. Status badges can use severity colors.
Validation can use problem colors. Primary commands can use accent color if
the Look-and-Feel supports it. Most form surfaces should remain quiet. A
business UI dominated by decorative color becomes harder to scan. The figure
screenshots for this chapter should show useful color, realistic spacing, and
clear hierarchy without turning the controls into a marketing page.

Professional custom controls also need responsive geometry. A unit field must
handle a longer unit abbreviation. A buddy field must handle a long path. A
split button must handle a localized command label. A table header must
handle translated column names and still leave room for sort/filter glyphs.
A calendar must handle six visible weeks. A token field must wrap or scroll
tokens intentionally. If a screenshot looks good only at one exact English
width, it is not a strong design demonstration.

Corusco helps aesthetics indirectly by reducing duplicated wiring. When text,
tooltip, help, and disabled reason come from resources and command metadata,
the visual component can focus on layout and state. When problems are
structured, decoration can be consistent. When table descriptors name columns,
header menus can be generated consistently. When component keys identify
widgets, screenshot tests can assert and document the intended visual states.
Maintainable code makes better visuals easier because state is not scattered.

\begin{principlebox}{Lessons Learned From Visual Design}
Professional appearance is mostly consistency under change. A widget that uses
UI defaults, stable metrics, real focus indicators, and screenshot-tested states
will age better than one whose colors and insets were chosen by eye in a single
theme.
\end{principlebox}

\section{Accessibility and Input Contracts}

\begin{roadmapbox}
This section makes keyboard and assistive-technology behavior part of the
component contract. It reviews names and descriptions, focus traversal,
shortcuts, popup invocation, non-color state, and the discipline of testing the
same operation through mouse and keyboard routes.
\end{roadmapbox}

Accessibility is not a separate pass at the end of component construction. It
is part of the input contract. If a user can click something, there should be
a keyboard and accessible path to the same operation unless the operation is
purely decorative. If a control communicates state visually, the same state
should be available as text or accessible description. If a problem can be
shown in a validation strip, the user should be able to navigate to the
problem target. These rules are easier to satisfy when the model is explicit.

For split buttons, the accessible contract should describe both surfaces. The
primary action has a name. The popup has a name. The shortcut to open the
popup should be discoverable. If the arrow is not a focus stop, the main
button description can mention Alt+Down or the chosen key. If the arrow is a
focus stop, it needs its own accessible name. Disabled reasons should be
available for both primary and secondary commands.

For buddy fields, the buddy button name should include the target field. "Clear
customer search" is better than "Clear." "Browse for report folder" is better
than "Browse." If the visible label labels only the text field, the wrapper
and buddy may still need descriptions. If validation applies to the whole
wrapper, the accessible description should expose the problem. Corusco field
keys and resource keys give the installer enough identity to generate stable
names.

For unit fields, the accessible value should include amount and unit. A screen
reader should not announce only the amount field and then a mysterious "kg
button." The unit segment should say that it changes the unit for the weight
field. If the unit is fixed, it should be part of the field description. If
the unit is editable, it should be reachable or invokable from the keyboard.
The problem message should distinguish amount and unit when that helps
recovery.

For table headers, accessible text should include sort state, filter state,
and popup availability where possible. A header that paints a sort glyph but
does not expose sort direction is incomplete. A filter icon that is color-only
is incomplete. A column visibility menu should use localized column names for
the user while retaining stable descriptor ids internally. F1 help should
route by help topic, not by header string.

For temporal selectors, accessible day names should include full dates and
state: selected, focused, today, disabled, outside current month, start, end,
or in range. A custom-painted calendar must work harder because Swing will
not automatically expose painted cells as children. If that is too much for a
small picker, prefer a grid of components. A professional calendar is not only
pretty; it is operable without a mouse.

Input maps should be documented near the component. A split button has popup
keys. A unit field has unit-popup keys. A calendar has day and month
navigation. A token field has Backspace and arrow policies. A validation strip
has activation and next-problem commands. A step panel has navigation
commands. These bindings should use \api{ActionMap}, not key listeners, so
tests and alternate input paths can invoke the same actions.

\begin{principlebox}{Lessons Learned From Accessibility Contracts}
Accessibility improves the ordinary engineering shape of a widget. Once a
control has a name, focus rule, keyboard command, state description, and
non-color feedback, it is also easier to test and easier for maintainers to
understand.
\end{principlebox}

\section{Failure Patterns in Practical Components}

\begin{roadmapbox}
This section collects the recurring ways practical components decay after a
demo: state hidden in renderers, listener-only workflows, display text used as
identity, stale async callbacks, popup leaks, and validation surfaces that do not
match the visible field.
\end{roadmapbox}

The most common failure pattern is text as identity. A split button decides
which command to run by reading button text. A table header stores widths by
visible header. A validation summary finds a field by label. A status badge
exports its localized label as a status code. These implementations work
until localization, copy editing, or user customization arrives. Corusco keys
exist to prevent this category of bug. Use them.

The second failure pattern is listener as model. A document listener parses
text, sets borders, enables Save, updates a tooltip, and starts a search. A
mouse listener opens a popup, computes menu items, writes selection, and logs
analytics. A calendar click handler changes visible month, selected date, and
committed value in one branch. These listeners are hard to test and easy to
bypass from keyboard input. Move durable decisions into models and commands.

The third failure pattern is child component leakage. A unit field exposes a
text field and button, so callers wire them differently on every screen. A
token field exposes its chip panel, so callers remove chips directly. A
wizard exposes its card layout, so callers jump pages without validation. A
custom component can expose child components for rare integration needs, but
its primary API should expose logical facts and commands.

The fourth failure pattern is visual-only state. A badge is red but has no
status code. A calendar paints disabled days but keyboard selection accepts
them. A table header paints a filter icon but the filter model is elsewhere
and can drift. A field has a red border but no structured problem. Visual
state should reflect model state. It should not be the only state.

The fifth failure pattern is missing close. Popup listeners remain attached.
Suggestion tasks complete after the view closes. Table-state controllers keep
writing preferences. Validation bindings update old components. Timers keep
running. In small demos this looks harmless. In business applications with
long-lived workspaces, it becomes memory leaks and stale UI updates. Every
installer should return a binding, subscription, controller, or scope entry.

The sixth failure pattern is test altitude mismatch. A screenshot is used to
prove parsing. A full UI test is used to prove duplicate-token policy. A unit
test is used to claim focus traversal. Use the lowest layer that proves the
claim, then add screenshots for visual integration. Corusco makes this easier
because fields, commands, problems, descriptors, and scopes can be tested
without rendering the whole screen.

\begin{principlebox}{Lessons Learned From Failure Patterns}
Most practical-component failures are not Swing limitations. They are missing
contracts. When identity, model ownership, event routing, cleanup, and problem
targets are explicit, ordinary Swing code is usually predictable enough for
business applications.
\end{principlebox}

\section{Corusco Implementation Map}

\begin{roadmapbox}
This section maps chapter concepts back to Corusco features. It identifies where
field keys, command keys, resource bundles, descriptors, problem models,
lifecycle scopes, and generated bindings reduce boilerplate without hiding the
Swing implementation from reviewers.
\end{roadmapbox}

It is useful to map each widget decision to a Corusco concept. Field identity
maps to \api{FieldKey}. Component lookup maps to \api{ComponentKey}. Command
identity maps to \api{ActionKey}. Visible text maps to \api{ResourceKey}.
Help maps to \api{HelpTopic}. Validation maps to \api{ProblemCode},
\api{ProblemTarget}, and \api{ProblemSet}. Table identity maps to
\api{TableKey} and \api{ColumnKey}. Lifecycle maps to \api{Subscription},
\api{Binding}, \api{BindingScope}, and \api{DetachableScope}. Task freshness
maps to \api{GenerationCounter}.

This map does not mean every custom component imports every Corusco type. A
reusable widget may remain plain Swing and expose a small interface. The
binding layer imports Corusco and connects keys, models, commands, and
problems. This keeps the reusable component from becoming framework-heavy
while still letting application code benefit from generated metadata. The
boundary is especially valuable for components that might later be replaced
by a third-party widget or custom-painted implementation.

For generated forms, custom fields can be treated as custom views over
generated field models. The processor may generate a field key and validation
metadata. The application writes an adapter from the generated field model to
the custom component. If the same custom component appears often, that adapter
becomes a reusable binding factory. The generator does not need to know every
visual control in advance. It needs to generate stable facts. The binding
layer consumes them.

For generated actions, split buttons and overflow bars can be ordinary
surfaces over generated command descriptors. The same action key can feed a
toolbar button, popup menu item, keyboard shortcut, command palette, and help
system. A split button chooses a primary action and secondary actions. An
overflow bar chooses visible and overflowed actions. Neither surface owns the
command. This is how the UI stays consistent when commands appear in several
places.

For generated tables, advanced headers can be ordinary installers over table
descriptors. The table descriptor supplies column keys, header resources,
tooltip resources, help topics, persistence ids, default widths, and
capabilities. The header installer supplies Swing behavior: mouse resolution,
keyboard popup action, renderer decoration, and menu construction. The table
state controller supplies persistence. Each part has a small job.

For problems and validation, custom widgets should consume structured
problems rather than strings. A unit field can show amount and unit problems.
A token field can show duplicate-token problems. A temporal selector can show
range-order problems. A validation strip can route to targets. The problem
model is independent of current component layout. That independence is what
lets the view change from plain fields to compound components without
rewriting validation logic.

For lifecycle, the construction code should be boring. Create a scope. Install
bindings. Add each binding to the scope. Close the scope when the view closes.
If an installer cannot return a closeable relationship, ask why. Some
relationships are genuinely permanent for the life of the application, but
screen-level custom components rarely are. Popups, document listeners,
validation observers, table state controllers, suggestion subscriptions, and
task handles all need cleanup owners.

This implementation map is also a migration map. A legacy screen can adopt
Corusco one layer at a time. First introduce keys and resources. Then wrap
commands. Then move validation into problem sets. Then bind custom fields to
field models. Then use table descriptors for headers and persistence. Then
close relationships through scopes. Each step reduces hidden state without
requiring every custom component to be rewritten at once.

\begin{principlebox}{Lessons Learned From the Corusco Map}
The maintainability gain comes from stable generated facts, not from removing
Swing from the design. Corusco should make custom widgets easier to wire,
localize, validate, test, and clean up while leaving their visual behavior
inspectable as ordinary Swing code.
\end{principlebox}

\section{Review Protocol for Practical Components}

\begin{roadmapbox}
This section gives reviewers a concrete protocol for evaluating a practical
component before it becomes part of a production screen. The path is: identify
the component contract, inspect model ownership, trace command and event paths,
check focus and accessibility, review Corusco integration, inspect visual states,
test lifecycle cleanup, and decide whether the component is ready to reuse.
\end{roadmapbox}

A practical component is rarely reviewed well by opening the screenshot first.
Screenshots are necessary, especially for compound controls whose purpose is to
look integrated, but a screenshot cannot reveal whether the selected unit is
part of the field value, whether a popup command has a stable key, whether a
table cell command uses the correct row after sorting, or whether a background
suggestion result can update a disposed screen. A review protocol gives the
team a repeatable order of questions so visual polish and behavioral correctness
support each other instead of competing.

The first review pass names the component boundary. The reviewer should be able
to point at the public view interface, model class, binding, or presenter method
and say what the application sees. If the only boundary is "this panel contains a
text field, a button, and a popup", the component is still too physical. The
reviewer should ask what value is edited, what commands are exposed, what
problems can occur, what focus targets exist, and what cleanup is required. The
answer does not need to be an elaborate framework. It does need to be explicit
enough that a test can address the same facts.

\begin{tabularx}{0.96\linewidth}{@{}p{0.24\linewidth}X@{}}
\toprule
Review pass & Question\\
\midrule
Boundary & What does the application see: value, commands, problems, focus, cleanup?\\
Model & Which durable facts are outside listeners and renderers?\\
Command & Which \api{Action} or Corusco command key owns each operation?\\
Events & How does a gesture become a model update or presenter command?\\
Focus & Can every operation be reached by keyboard in a predictable order?\\
Visual states & Are normal, hover, focused, disabled, invalid, and loading states shown?\\
Lifecycle & Who closes listeners, popups, timers, subscriptions, and async tasks?\\
\bottomrule
\end{tabularx}

The second pass inspects model ownership. Every practical widget in this chapter
has a fact that is tempting to hide. Split buttons hide the default command.
Buddy fields hide the relationship between the text and the button. Unit fields
hide the unit in a visual segment. Table headers hide column identity in header
text. Tag clouds hide chip identity in labels. Interactive cells hide commands
in renderer geometry. Temporal selectors hide draft dates inside popup widgets.
The review question is always the same: if the screen repaints, if the popup is
rebuilt, if the table sorts, or if the Look-and-Feel changes, where does that
fact still live?

\begin{lstlisting}[caption={A review can ask for the component contract before the layout.}]
interface PracticalComponentReview {
    Object valueContract();
    List<Action> exposedActions();
    ProblemTarget problemTarget();
    List<FocusTarget> focusTargets();
    AutoCloseable lifecycleHandle();
}
\end{lstlisting}

This interface is not meant to be implemented literally. It is a review aid. If
the team cannot answer one of those methods in prose, the component will be hard
to maintain. A pure status badge may not edit a value, but it still has status
identity and accessibility text. A table cell renderer may not have its own
focus target, but the table cell selection and keyboard command are focus
contracts. A calendar popup may not return an \api{AutoCloseable} directly, but
some binding or scope must own its listeners.

The third pass traces one command from resource key to effect. Choose one
operation such as Export PDF, Clear field, Remove tag, Change unit, Hide column,
Approve row, or Apply date range. Find the generated command key or action id.
Find the resource text and help. Find enabled state and disabled reason. Find the
keyboard binding. Find the mouse binding. Find the presenter or model method
that performs the effect. Find the refresh that makes the component show the new
state. This trace should be short. If it crosses anonymous listeners that cannot
be named, the command is not reviewable enough.

\begin{detailbox}{Command Trace Heuristic}
For any visible operation, the reviewer should be able to answer: "what command
is this?", "why is it enabled?", "how can I invoke it without a mouse?", "what
model state changes?", and "what closes when the screen closes?"
\end{detailbox}

The fourth pass checks coordinate translations. This pass is essential for
tables, headers, and interactive cells, but it is useful elsewhere too. Table
events arrive in view coordinates and must be converted to model identity before
business logic runs. Header clicks arrive on a visible column and must be mapped
to a descriptor key. Tag-chip hit testing may happen in wrapped visual
coordinates and must become a tag id. Calendar clicks happen in grid coordinates
and must become a \api{LocalDate} under the current visible month and zone
policy. Reviewers should search for places where the implementation stores a
view index beyond the current event callback.

The fifth pass checks focus policy. The review should not merely tab through the
component once. It should ask what the intended focus stops are and compare the
implementation to that design. A unit field may be one focus stop with Alt+Down
for units. A buddy field may include the Browse button as a separate focus stop
but exclude Clear. A tag cloud may keep focus in the editor while moving logical
focus among chips. A table may keep keyboard focus on the table while invoking
cell commands. Each policy is acceptable when it is intentional and documented
by input-map tests.

\begin{tabularx}{0.96\linewidth}{@{}p{0.24\linewidth}X@{}}
\toprule
Focus check & Evidence\\
\midrule
Entry & Tab lands on the intended first target.\\
Exit & Tab and Shift+Tab leave by predictable paths.\\
Popup & Alt+Down, Space, Enter, Escape, and Shift+F10 work where expected.\\
Removal & Deleting a chip, row command, or field value leaves focus somewhere useful.\\
Validation & Focus can move to the field or cell targeted by a problem.\\
Disposal & Closing a popup or screen does not leave focus on a disposed component.\\
\bottomrule
\end{tabularx}

The sixth pass checks accessibility and language. This is not a separate
late-stage audit; it is a way to reveal unclear component contracts. If the team
cannot write an accessible name for a component, the component probably does not
have a clear purpose. If a status can only be described as "the red one", the
status model is insufficient. If a table cell command cannot be reached without
clicking a painted icon, the interaction is incomplete. If a field error appears
only as a border color, the problem model is not communicating enough.

The seventh pass checks Corusco integration. Generated keys should appear where
the application meaning appears. Field keys should target values and problems.
Command keys should target operations. Component keys should help tests locate
important surfaces without depending on English labels. Table descriptors should
identify columns. Resource keys should supply visible text and accessible text.
Lifecycle scopes should own subscriptions. The review should be suspicious of
code that reconstructs any of those facts from text, index position, or current
child order.

\begin{lstlisting}[caption={A review-friendly binding keeps generated identity visible.}]
final class ReviewableBinding implements Binding {
    private final FieldKey<?> fieldKey;
    private final CommandKey clearCommand;
    private final ComponentKey fieldComponent;
    private final Subscription modelSubscription;

    @Override
    public void close() {
        modelSubscription.close();
    }
}
\end{lstlisting}

The eighth pass checks visual states. A component that looks attractive only in
the happy path is not finished. Review normal, focused, hover, pressed, disabled,
invalid, loading, empty, long-text, high-DPI, large-font, and narrow-width
states. For tables, also review selected row, focused cell, sorted column,
filtered row count, and disabled command states. For temporal controls, review
today, selected date, out-of-range dates, disabled dates, other-month dates, and
invalid spans. For tag clouds, review overflow, wrapped rows, focused chip,
invalid chip, loading suggestions, and failed suggestions.

Visual review should use real Swing demo screenshots from the book's figure
pipeline. Hand-generated mockups can be useful for ideation, but chapter figures
that teach Swing and Corusco techniques should demonstrate actual Swing
components wherever practical. That keeps the screenshot honest: fonts, borders,
row heights, insets, antialiasing, and Look-and-Feel defaults come from the same
environment the code uses. A beautiful mock that cannot be built in Swing is a
design debt, not a documentation asset.

The ninth pass checks lifecycle. Practical components create small resources:
document listeners, property-change listeners, row-sorter listeners, popup-menu
listeners, timers, background requests, model subscriptions, and generated
bindings. The reviewer should ask where each is closed. This is especially
important for popup controls and async suggestions because the leak may not show
until a screen is opened and closed many times. A Corusco scope or binding
lifecycle makes the answer explicit.

\begin{tabularx}{0.96\linewidth}{@{}p{0.24\linewidth}X@{}}
\toprule
Resource & Cleanup owner\\
\midrule
Document listener & Field binding or component scope.\\
Popup listener & Popup installer or command-surface binding.\\
Sorter listener & Table binding that installed header behavior.\\
Async suggestion & Search controller tied to screen lifecycle.\\
Timer & Component scope, stopped on dispose.\\
Problem subscription & Validation binding.\\
\bottomrule
\end{tabularx}

The tenth pass decides reuse readiness. A prototype can be screen-local if it is
honest about that status. A reusable component needs a cleaner public contract,
resource strategy, test surface, and lifecycle story. Do not promote a widget to
a shared package because it looks useful in two places. Promote it when the
second usage proves that the model and command contracts are stable enough.
Corusco can reduce the cost of reuse by making field, command, resource, and
problem identities uniform across screens.

Reviewers should also know when not to generalize. A table cell command for one
screen may remain a screen-specific column policy. A date span selector used by
one reporting module may remain a reporting component. A general library should
appear only when variation is understood. The protocol helps here because it
separates generic facts from screen-specific facts. Generic facts belong in the
component. Screen-specific facts belong in descriptors, bindings, presenters, or
command contributors.

\begin{principlebox}{Lessons Learned From the Review Protocol}
A practical component is ready when a reviewer can trace value, command, focus,
problem, visual state, and lifecycle without reading accidental child-component
details. The protocol makes aesthetics and maintainability reinforce each other:
screenshots prove the Swing surface, while Corusco keys and models prove that
the surface is attached to stable application meaning.
\end{principlebox}

\section{Testing and Review}

\begin{roadmapbox}
This final section turns the construction advice into verification practice. It
covers model tests, event-path tests, focus and keyboard tests, popup lifecycle
tests, accessibility smoke checks, screenshot tests, and acceptance rules for
each major component family in the chapter.
\end{roadmapbox}

Compound components need a testing pyramid that matches their contracts. Model tests are
fast and prove value rules. Swing interaction tests prove actions, focus, keyboard, and
popup paths. Screenshot tests prove visual integration across Look-and-Feel, density,
font, and high-DPI assumptions. Accessibility smoke checks prove that names and
descriptions exist and that state is not color-only. Lifecycle tests prove listeners and
timers stop when the owner closes.

\begin{tabularx}{0.96\linewidth}{@{}p{0.28\linewidth}X@{}}
\toprule
Test type & What it should prove\\
\midrule
Model test & Parsing, formatting, selection, validation, enabled state, and commit or cancel rules.\\
Action test & Primary and secondary commands invoke the right presenter method and observe enabled state.\\
Keyboard test & Every mouse-only-looking affordance has a named keyboard path.\\
Popup test & Menus open from mouse and keyboard, refresh state, and do not mutate unrelated selection.\\
Screenshot test & Borders, baselines, density, disabled state, problem state, and popup affordances look integrated.\\
Lifecycle test & Listeners, timers, subscriptions, and pending tasks are closed with the view.\\
\bottomrule
\end{tabularx}

A screenshot test should not be the first test. It is tempting because custom controls
are visual. But a screenshot cannot explain whether a unit change converted the amount
correctly, whether a calendar range uses exclusive end, or whether a split button's
arrow remains enabled when the primary action is disabled. Those are model and action
tests. Use screenshots for what only pixels can prove: alignment, clipping, colors,
borders, focus rings, popup placement, and density.

\begin{lstlisting}[caption={A model test can avoid Swing entirely.}]
@Test
void unitChangeKeepsDraftTogether() {
    MeasurementFieldModel model =
            new MeasurementFieldModel(MeasurementUnit.KILOGRAM);

    model.setRawAmount("120");
    model.setUnit(MeasurementUnit.POUND);

    assertThat(model.draft().rawAmount()).isEqualTo("120");
    assertThat(model.draft().unit()).isEqualTo(MeasurementUnit.POUND);
    assertThat(model.draft().canCommit()).isTrue();
}
\end{lstlisting}

The exact test framework is not important here. The important shape is that the unit
field's core behavior can be tested before any Swing component is created. If the only
way to test a unit change is to click a tiny arrow in a rendered window, too much state
is trapped in the component.

Keyboard tests should name the action rather than rely only on physical key delivery
when possible. It is useful to test that Alt+Down is mapped to
\api{openPopup}; it is also useful to invoke \api{openPopup} directly and
assert menu state. Physical key tests can be fragile across platforms. Action map tests
are stable and verify the component contract.

Popup lifecycle tests should catch stale menus. Open the menu for column A, close it,
open for column B, and verify the actions target B. Disable a secondary split-button
action, open the menu, and verify the item is disabled. Close the view, update the
command model, and verify the old component does not receive updates. These tests
prevent subtle leaks that are hard to see in screenshots.

Focus traversal tests should run on the EDT and inspect the actual focus cycle where
feasible. A buddy field that should be one focus stop should not add its clear button to
the normal Tab order. A browse button that should be reachable should appear in the
cycle. A date-time span selector should move predictably through start date, start time,
end date, end time, presets, and buttons. These tests are small but they protect
keyboard users and fast data entry workflows.

Accessibility smoke checks can be pragmatic. Verify that custom controls and buddy
buttons have nonblank accessible names. Verify that state descriptions include selected
unit, sort direction, severity, or disabled reason where appropriate. Verify that
color-only badges also have text. Full accessibility testing requires real assistive
technology, but basic metadata checks catch many avoidable omissions.

Acceptance gates should be written per component family. A generic statement such as
"the custom control is tested" is too vague. A split button needs command tests and
popup tests. A unit field needs parsing, formatting, and unit-change tests. A table
header needs coordinate translation and descriptor identity tests. A tag cloud needs
keyboard navigation and stale-suggestion tests. Interactive cells need hit-zone and
row-id tests. Temporal panels need date arithmetic and commit or cancel tests. These
gates make the review repeatable and help future maintainers know which behavior was
intentional.

\begin{tabularx}{0.96\linewidth}{@{}p{0.25\linewidth}X@{}}
\toprule
Component family & Minimum acceptance gates\\
\midrule
Split button & Primary command, secondary popup, disabled reason, keyboard popup, cleanup.\\
Buddy field & Document sync, command action, focus policy, problem decoration, chooser cancel.\\
Unit field & Raw text, parsed amount, selected unit, conversion policy, problem targeting.\\
Table header & View/model conversion, descriptor key, sort glyph, popup routing, persistence.\\
Temporal picker & Draft value, commit value, bounds, disabled days, zone, endpoint policy.\\
Tag cloud & Stable tag id, query state, suggestions, chip focus, removal, async cancellation.\\
Interactive cell & Renderer/editor split, hit zones, row id, keyboard command, stale popup.\\
\bottomrule
\end{tabularx}

The test names should tell the design story. A test named
\texttt{buttonWorks()} does not help a maintainer. A test named
\texttt{arrowRemainsEnabledWhenPrimaryCommandIsDisabled()} explains an important
split-button rule. A test named
\texttt{staleSuggestionsDoNotReplaceNewerTagQuery()} explains async policy. A
test named \texttt{cellCommandUsesRowIdAfterSorterChanges()} explains why
coordinate conversion exists. Good test names are a cheap form of documentation
for component decisions that are otherwise easy to simplify incorrectly.

\begin{lstlisting}[caption={Test names should preserve design decisions.}]
@Test
void arrowRemainsEnabledWhenPrimaryCommandIsDisabled() {
    commands.exportPdf().setEnabled(false);
    commands.exportCsv().setEnabled(true);

    splitButton.refresh(commands.snapshot());

    assertThat(splitButton.primaryEnabled()).isFalse();
    assertThat(splitButton.popupEnabled()).isTrue();
}

@Test
void cellCommandUsesRowIdAfterSorterChanges() {
    table.sortBy(OrderColumns.AMOUNT);

    invokeActionAtSelectedCell(OrderColumns.APPROVAL);

    assertThat(commands.lastInvocation().rowId()).isEqualTo(expectedOrderId);
}
\end{lstlisting}

Screenshot tests need stable scenarios. Each component family should have at least one
normal scenario and one stress scenario. Normal shows the control as it should appear in
the book. Stress shows long labels, disabled commands, problem state, or large fonts.
The stress shot may not be as pretty, but it catches the defects that reach production:
clipped unit labels, buddy buttons that steal all width, tag chips that overlap, table
header menus that cover the wrong cell, and date span panels whose buttons move under
larger fonts.

When screenshot tests fail, do not approve them by habit. A small pixel difference can
be harmless after a Look-and-Feel upgrade, but it can also reveal a real metrics change.
The review should ask whether the component still communicates the same hierarchy,
state, and affordances. If the failure is harmless, update the baseline with a note.
If the failure changes focus rings, disabled contrast, error placement, or popup
alignment, treat it as a component regression.

Tests should also exercise cleanup. A practical pattern is to create the component,
connect it to a model, perform one interaction, close the binding or scope, then mutate
the model again. The disposed component should not receive the update. For async
components, start a request, close the scope, complete the request, and verify that no
popup opens and no view method is called. These tests are not glamorous, but they catch
the slow leaks that make long-running desktop applications feel unreliable.

\begin{lstlisting}[caption={Lifecycle tests mutate the model after the view closes.}]
@Test
void closedBindingDoesNotReceiveModelUpdates() {
    TagCloudBinding binding = bindTagCloud(view, model);

    binding.close();
    model.addTag(new TagId("late"));

    assertThat(view.refreshCount()).isEqualTo(1);
}
\end{lstlisting}

Finally, review the documentation example itself. A code sketch in the book should have
the same design shape as recommended production code even when it omits imports or
surrounding infrastructure. If the prose says commands are named, the listing should
not rely on anonymous listener-only state. If the prose says descriptor keys matter,
the listing should not derive behavior from header text. If the prose says Corusco
keeps bindings maintainable, the example should show keys, descriptors, problem
targets, resources, or lifecycle scopes doing real work.

Review should start with ownership. For each custom control, ask where the value lives,
where commands live, where disabled state lives, where problems live, where resources
live, where help lives, and where cleanup lives. Then walk input paths: mouse, keyboard,
focus, popup, cancel, and close. Finally walk visual states: normal, hover, pressed,
focused, disabled, error, high contrast, large font, and narrow width. A component that
survives this review is usually ready to reuse.

\begin{principlebox}{Reusable Means Reviewable}
A reusable Swing component is not one that can be dropped onto any panel. It is one
whose value, commands, focus policy, resources, problems, and lifecycle can be
understood without reading every mouse listener.
\end{principlebox}

\begin{exercisebox}
\begin{enumerate}
\item Build a split button for Export. Make the primary action export PDF and
the popup offer CSV and XLSX. Verify that the popup remains available when PDF export is
disabled.
\item Replace a text field plus separate unit combo box with one measurement
field wrapper. Write model tests for invalid amount text and unit changes.
\item Add a header context menu to a descriptor-backed table. The menu should
hide columns by descriptor id and route help by help topic.
\item Design a date-time span selector model for a report filter. Specify
whether the end instant is inclusive or exclusive and test the boundary.
\item Choose one status badge in an existing screen. Make sure text,
tooltip, accessible description, and exported value all come from the same status model.
\item Write a screenshot test for a buddy field at a larger font size and a
model test for the same field's clear action.
\end{enumerate}
\end{exercisebox}

The practical lesson is deliberately conservative. Build compound components from
ordinary Swing when ordinary Swing gives you focus, text editing, accessibility, and
Look-and-Feel behavior for free. Introduce custom painting only when the visual contract
really needs it. Keep models explicit. Keep commands named. Keep validation structured.
Keep headers descriptor-backed. Keep temporal policy visible. Then the business
application can have richer controls without turning every form into a private widget
framework.

\begin{principlebox}{Lessons Learned From Testing and Review}
Tests are useful when they preserve the component decisions, not merely when they
exercise the rendered pixels. Model tests protect value semantics, action tests
protect command ownership, interaction tests protect keyboard and popup paths,
screenshot tests protect visual integration, and lifecycle tests protect long
running screens. Together they make rich Swing components safe to reuse.
\end{principlebox}
